\documentclass[11pt,a4paper]{amsart}

\usepackage[utf8]{inputenc}
\usepackage[T1]{fontenc}
\usepackage{lmodern}
\usepackage{amsmath,amssymb,amsthm,mathtools}
\usepackage{enumitem}
\usepackage{microtype}
\usepackage[hidelinks]{hyperref}
\usepackage[margin=1.1in]{geometry}

\numberwithin{equation}{section}

\newtheorem{theorem}{Theorem}[section]
\newtheorem{proposition}[theorem]{Proposition}
\newtheorem{lemma}[theorem]{Lemma}
\newtheorem{corollary}[theorem]{Corollary}

\theoremstyle{definition}
\newtheorem{definition}[theorem]{Definition}
\newtheorem{remark}[theorem]{Remark}
\newtheorem{example}[theorem]{Example}

\newcommand{\R}{\mathbb{R}}
\newcommand{\C}{\mathbb{C}}
\newcommand{\Z}{\mathbb{Z}}
\newcommand{\Q}{\mathbb{Q}}
\newcommand{\F}{\mathbb{F}}
\newcommand{\Kcal}{\mathcal{K}}
\newcommand{\Pcal}{\mathcal{P}}
\newcommand{\Hcal}{\mathcal{H}}
\newcommand{\Dcal}{\mathcal{D}}
\newcommand{\ind}{\operatorname{ind}}
\newcommand{\wind}{\operatorname{wind}}
\newcommand{\Lip}{\operatorname{Lip}}
\newcommand{\inte}{\operatorname{int}}
\newcommand{\vol}{\operatorname{vol}}
\newcommand{\one}{\mathbf{1}}
\newcommand{\avg}[2]{\langle #1 \rangle_{#2}}

\title[Counterexamples to the Generalized NRR Problem]{Non-Prime-Power Counterexamples to the Generalized Nandakumar--Ramana Rao Problem in Dimension at Least Three}

\author{GPT 6.2 Astra}
\author{GPT 6.1 Sol}
\author{Gemini 4 Neon}
\author{Claude Fable 5.2}
\date{November 27, 2026}

\begin{document}

\begin{abstract}
Let $d \ge 3$ and let $n \ge 2$ be an integer that is not a prime power. We prove that for every prescribed full-dimensional compact convex body $K \subset \R^d$ and every prescribed probability measure $\mu$ on $K$ that is absolutely continuous with respect to Lebesgue measure, there exist $d-1$ Hausdorff-continuous real-valued functionals $\varphi_1, \dots, \varphi_{d-1}$ defined on all nonempty compact convex subsets of $\R^d$ such that no convex $n$-partition of $K$ into pieces of equal $\mu$-measure simultaneously equalizes all $\varphi_k$. Moreover, the same functionals obstruct approximate equipartitions with a uniform positive gap. Combined with the known existence theorem for prime powers due to Karasev, Hubard, and Aronov and to Blagojevi\'c and Ziegler, this yields a complete pointwise classification of the generalized Nandakumar--Ramana Rao problem in dimensions $d \ge 3$: universal fairness holds for $(K,\mu)$ if and only if $n$ is a prime power.

For the uniform $d$-dimensional cube $K = [-1,1]^d$, we further prove a dimension-uniform polynomial strengthening: for each non-prime-power $n \ge 2$, there exists a single pair of four-variable polynomials $P_n, Q_n \in \Q[X_1,X_2,X_3,X_4]$ such that, in every dimension $d \ge 3$, the two polynomial combinations of four low-order volume moments together with the $d-3$ remaining coordinate first moments obstruct all equal-volume convex $n$-partitions of $[-1,1]^d$.
\end{abstract}

\keywords{Nandakumar--Ramana Rao problem, convex equipartitions, prime powers, winding number, Brunn--Minkowski inequality}
\subjclass[2020]{Primary 52A37; Secondary 52A20, 52A40, 55M20}

\maketitle

\section{Introduction}
\label{sec:intro}

The Nandakumar--Ramana Rao (NRR) problem~\cite{NRR2012} originally asked whether every convex polygon in the plane can be partitioned into $n$ convex pieces of equal area and equal perimeter. In higher dimensions, the natural generalized problem replaces area and perimeter by an absolutely continuous probability measure $\mu$ on a convex body $K \subset \R^d$ and $d-1$ continuous real-valued functionals on convex sets.

Let $\Kcal(\R^d)$ denote the space of all nonempty compact convex subsets of $\R^d$, equipped with the Hausdorff metric. By a \emph{convex body} we mean a full-dimensional element $K \in \Kcal(\R^d)$ (that is, $\inte K \ne \emptyset$). A \emph{convex $n$-partition} of $K$ is a family $C_1, \dots, C_n$ of convex bodies in $\R^d$ with mutually disjoint interiors such that $\bigcup_{i=1}^n C_i = K$. Given a probability measure $\mu$ on $\R^d$ supported on $K$ (i.e.\ $\mu(K) = 1$) and absolutely continuous with respect to $d$-dimensional Lebesgue measure, we say that a convex $n$-partition $\{C_i\}_{i=1}^n$ is \emph{$\mu$-equipartitioned} if $\mu(C_i) = 1/n$ for all $i = 1, \dots, n$, and \emph{fair} for a tuple of functionals $\varphi = (\varphi_1, \dots, \varphi_{d-1}) : \Kcal(\R^d) \to \R^{d-1}$ if in addition $\varphi(C_1) = \dots = \varphi(C_n)$.

For prime powers $n = p^a$, Karasev, Hubard, and Aronov~\cite{KHA2014} and Blagojevi\'c and Ziegler~\cite{BZ2014} proved using equivariant configuration-space obstructions and power diagrams that for every $d \ge 2$, every convex body $K \subset \R^d$, every absolutely continuous probability measure $\mu$ on $K$, and every continuous $\varphi : \Kcal(\R^d) \to \R^{d-1}$, a fair convex $n$-partition always exists. In dimension $d = 2$, Akopyan, Avvakumov, and Karasev~\cite{AAK2018} established existence for \emph{all} integers $n \ge 2$ (for one absolutely continuous measure and one continuous functional; see~\cite[Theorem~7.1]{AAK2018}).

In dimensions $d \ge 3$, whether universal fairness persists for non-prime-power $n$ (starting with $d=3$ and $n=6$) remained open and was explicitly formulated as a conjecture in~\cite[Conjecture~1.2]{BKS2023}. The main purpose of this paper is to disprove this conjecture in the strongest pointwise sense: not only do counterexamples exist for every $d \ge 3$ and every non-prime-power $n$, they exist for \emph{every prescribed} convex body $K$ and \emph{every prescribed} absolutely continuous probability measure $\mu$.

\begin{theorem}[Prescribed-body, prescribed-measure counterexamples]
\label{thm:prescribed}
Let $d \ge 3$ and let $n \ge 2$ be an integer that is not a prime power. Let $K \subset \R^d$ be a prescribed convex body, and let $\mu$ be a prescribed probability measure on $\R^d$ that is absolutely continuous with respect to Lebesgue measure and satisfies $\mu(K) = 1$. Then there exist $d-1$ Hausdorff-continuous functionals $\varphi_1, \dots, \varphi_{d-1} : \Kcal(\R^d) \to \R$ such that for every convex $n$-partition $C_1, \dots, C_n$ of $K$ with $\mu(C_1) = \dots = \mu(C_n) = 1/n$, there exist $k \in \{1,\dots,d-1\}$ and $i, j \in \{1,\dots,n\}$ with $\varphi_k(C_i) \ne \varphi_k(C_j)$.
\end{theorem}

The functionals in Theorem~\ref{thm:prescribed} depend on $d, n, K, \mu$ and are permitted to depend on position and orientation (they are not required to be rigid-motion invariant). In particular, Theorem~\ref{thm:prescribed} applies directly to the unit Euclidean ball $B^d$ with its normalized Lebesgue measure, as well as to partitions obtained by intersecting a convex partition of the ambient space $\R^d$ with $K$. If $\mu$ is any finite absolutely continuous measure with $\mu(K) > 0$, one simply restricts and renormalizes $\mu$ to $K$.

By compactness of the space of convex $n$-partitions, exact obstruction immediately upgrades to robust obstruction of approximate equipartitions.

\begin{corollary}[Robust counterexamples]
\label{cor:robust}
Under the hypotheses of Theorem~\ref{thm:prescribed}, the same continuous functionals $\varphi_1, \dots, \varphi_{d-1} : \Kcal(\R^d) \to \R$ satisfy the following stability property: there exists $\varepsilon > 0$ such that for every convex $n$-partition $C_1, \dots, C_n$ of $K$ satisfying $|\mu(C_i) - 1/n| < \varepsilon$ for all $i$, there exist $k, i, j$ with $|\varphi_k(C_i) - \varphi_k(C_j)| > \varepsilon$.
\end{corollary}

Combining Theorem~\ref{thm:prescribed} with the prime-power existence theorem of~\cite{KHA2014,BZ2014} gives a complete pointwise classification in dimensions $d \ge 3$.

\begin{corollary}[Pointwise classification for $d \ge 3$]
\label{cor:classification}
Let $d \ge 3$ and $n \ge 2$. For any convex body $K \subset \R^d$ and any absolutely continuous probability measure $\mu$ on $\R^d$ with $\mu(K) = 1$, every $(d-1)$-tuple of continuous functionals on $\Kcal(\R^d)$ admits a fair convex $n$-partition of $(K,\mu)$ if and only if $n$ is a prime power.
\end{corollary}

Our second main result is an explicit, dimension-uniform polynomial strengthening for the uniform $d$-cube $K = [-1,1]^d$. Write a point of $\R^d$ as $(x_1, \dots, x_{d-1}, z)$ with $z = x_d$, and for any integrable function $g$ and measurable set $C \subset \R^d$, write the normalized volume integral
\[
\avg{g}{C} := \frac{n}{2^d} \int_C g(x)\,dx.
\]
Define the four-moment map $M : \Kcal(\R^d) \to \R^4$ by
\begin{equation}
\label{eq:four-moments}
M(C) := \bigl(\avg{x_1}{C},\; \avg{x_2}{C},\; \avg{(1+x_1)z}{C},\; \avg{x_2 z}{C}\bigr).
\end{equation}

\begin{theorem}[Dimension-uniform rational polynomial obstruction on cubes]
\label{thm:cube}
Let $n \ge 2$ be an integer that is not a prime power. There exist two rational polynomials $P_n, Q_n \in \Q[X_1,X_2,X_3,X_4]$ (depending only on $n$, independent of $d$) such that for every dimension $d \ge 3$, the $d-1$ continuous functionals on $\Kcal(\R^d)$ given by
\[
C \longmapsto P_n(M(C)), \qquad C \longmapsto Q_n(M(C)), \qquad C \longmapsto \avg{x_j}{C} \quad (3 \le j \le d-1)
\]
admit no fair equal-volume convex $n$-partition of the cube $[-1,1]^d$. In dimension $d = 3$, there are only the first two functionals $P_n(M(C))$ and $Q_n(M(C))$, and moreover there exists $\gamma_n > 0$ such that every equal-volume convex $n$-partition of $[-1,1]^3$ has two pieces differing by more than $\gamma_n$ in either $P_n(M(\cdot))$ or $Q_n(M(\cdot))$.
\end{theorem}

The four integrands in~\eqref{eq:four-moments} have degree at most $2$, and the outer polynomials $P_n, Q_n$ have rational coefficients. Our proof establishes their existence via rational Stone--Weierstrass approximation and does not extract explicit degree bounds or explicit coefficients.

\subsection*{Relation to recent work and method of proof}

In a contemporaneous manuscript~\cite{Akopyan2026}, Akopyan constructs counterexamples for $n = 6$ in dimensions $d \ge 24$ on the unit ball (and $d \ge 18$ in the accompanying repository~\cite{AkopyanRepo}) using sixfold Tverberg maps and finitely many marked points, with existential quantification over the measure (sufficiently concentrated measures) and functionals. Our approach is completely different in both mechanism and scope: it covers all $d \ge 3$ and all non-prime-power $n \ge 2$, quantifies universally over every prescribed convex body $K$ and every prescribed absolutely continuous probability measure $\mu$, and builds single-cell perturbations directly from the geometry of two adjacent interfaces in a horizontal layering.

The proof rests on four ingredients:
\begin{enumerate}[label=(\arabic*)]
\item \textbf{A planar local obstruction (Section~\ref{sec:planar}).} Because $\gcd_{0<j<n}\binom{n}{j} = 1$ precisely when $n$ is not a prime power, we build a positively homogeneous homeomorphism $h : \C \to \C$ satisfying the exact antipodal identity $h(z) - h(-z) = 2z$ whose branch sums $g_j(u) = j h(u) + (n-j) h(-u)$ have prescribed winding numbers $d_j$ with $\frac12 \sum_{j=0}^n \binom{n}{j} d_j = 0$. Passing to a connected compact oriented Riemann surface $S_T$ of solutions to $z_i^2 = q + 4^{j(i)}$ ($i \ne 2$), the total boundary winding number of the residual map vanishes; hence the residual extends without zeros across $S_T$ and descends through a proper map to a bounded continuous single-cell correction $D(p_1, -p_2)$ depending \emph{only on the two interfaces bordering the second layer}.
\item \textbf{Cylinder rigidity and localization (Sections~\ref{sec:cylinder}--\ref{sec:ball-localization}).} Inside a symmetric cylinder $A = D \times [-h,h]$, a convex partition whose pieces all have zero horizontal first moment must consist purely of horizontal slabs. Combining an interior cylinder integral with a small $\delta$-weighted global integral forces every zero-moment partition into an arbitrarily small neighborhood of the compact family $\Hcal$ of horizontal quantile partitions.
\item \textbf{Uniform local interface charts and global assembly (Sections~\ref{sec:ball-assembly}--\ref{sec:general-body}).} Near $\Hcal$, every convex partition is uniquely ordered from bottom to top by small-slope affine interfaces, and the second piece determines both of its bounding interfaces from its own shape alone. Rescaling the planar correction by $\sigma^2$ and extending globally by Tietze's theorem preserves a positive gap outside the neighborhood while applying the planar obstruction inside. For a general convex body $K$, a maximal inscribed ball produces an interior chord between parallel supporting hyperplanes, and an explicit positive density $g \ge 1$ rebalances the horizontal first moment of $K$ to zero.
\item \textbf{Four-moment inversion on cubes and Brunn--Minkowski product rigidity (Sections~\ref{sec:cube-3d}--\ref{sec:cube-highd}).} On $[-1,1]^3$, the four moments~\eqref{eq:four-moments} algebraically invert the two bounding slopes of the second layer in a neighborhood of horizontal slabs, allowing the global perturbation to be written directly as a continuous function on $\R^4$ and approximated by rational polynomials $P_n, Q_n$. In dimensions $d > 3$, vanishing of the extra coordinate first moments forces every equal-volume candidate partition to be a strict product $B_i(0) \times [-1,1]^{d-3}$ via the concavity of $\vol_3(B_i(r))^{1/3}$, the equality case of the three-dimensional Brunn--Minkowski inequality, and a second-moment cancellation identity; hence the exact same pair $(P_n, Q_n)$ applies in every dimension $d \ge 3$.
\end{enumerate}

\medskip
\noindent\textbf{Exact arithmetic companion script.}
A self-contained Python script \texttt{verify.py} accompanies the repository for this paper and checks by exact rational arithmetic the internal binomial gcd criterion for $n = 2, \dots, 200$, the piecewise-linear winding realization and end-winding sums for $n \in \{6, 10, 12, 15\}$, all $324$ exact interpolation nodes for the four-moment inversion formulas on the cube, the $3$-dimensional ball cap bound at $1/\sqrt{2}$, and the positive-weight rebalancing identity on an asymmetric pyramid. The accompanying Lean~4 project proves both directions of Corollary~\ref{cor:classification} for each prescribed body and measure. Its final theorem is \texttt{NRR.fullClassification}.

% ============================================================================
\section{A Planar Local Obstruction for Non-Prime-Powers}
\label{sec:planar}
% ============================================================================

Throughout this section we identify $\R^2$ with $\C$, taking counterclockwise orientation as positive for winding numbers. Positive homogeneity of a map $\C \to \C$ always means homogeneity under multiplication by positive real scalars.

\begin{proposition}[Planar local lemma]
\label{prop:local-lemma}
Let $n \ge 2$ be an integer that is not a prime power. There exist a continuous positively homogeneous map $f : \C \to \C$ of degree $2$ (that is, $f(t w) = t^2 f(w)$ for all $t > 0$ and $w \in \C$), bounded continuous functions $\Delta_i : \C \to \C$ for $i \in \{1,\dots,n\} \setminus \{2\}$, and a bounded continuous function $D : \C^2 \to \C$ such that:
\begin{enumerate}[label=\textup{(\roman*)}]
\item \emph{(Balanced rigidity)} If $w_1, \dots, w_n \in \C$ satisfy $\sum_{i=1}^n w_i = 0$ and $f(w_1) = \dots = f(w_n)$, then $w_1 = \dots = w_n = 0$.
\item \emph{(Local obstruction)} For every tuple $p_0, p_1, \dots, p_n \in \C$ with $p_0 = p_n = 0$, setting $w_i := p_i - p_{i-1}$ for $1 \le i \le n$, the $n$ complex numbers
\[
f(w_2) + D(p_1, -p_2), \qquad f(w_i) + \Delta_i(w_i) \quad (i \ne 2)
\]
cannot all be equal.
\end{enumerate}
\end{proposition}

Crucially, the correction $D(p_1, -p_2)$ in the second term depends only on the two adjacent interface parameters $p_1$ and $p_2$. This is what allows the obstruction to be implemented by single-cell functionals on convex partitions.

\subsection{Realizing prescribed winding numbers with an exact antipodal identity}

\begin{lemma}
\label{lem:binom-gcd}
Let $n \ge 2$. Then
\[
\gcd\left\{\binom{n}{j} : 1 \le j \le n-1\right\} =
\begin{cases}
p & \text{if } n = p^a \text{ for a prime } p \text{ and } a \ge 1, \\
1 & \text{if } n \text{ is not a prime power.}
\end{cases}
\]
\end{lemma}

\begin{proof}
Fix any prime $p$ and write $n = p^a b$ with $a \ge 0$ and $p \nmid b$. In $\F_p[X]$, the Frobenius endomorphism gives $(1+X)^n = (1+X^{p^a})^b$. If $b > 1$ and $a \ge 1$, then $1 \le p^a < n$ and the coefficient of $X^{p^a}$ is $b \not\equiv 0 \pmod{p}$; if $a = 0$, then $p \nmid n = \binom{n}{1}$. Thus if $n$ is not a prime power, no prime divides all internal binomial coefficients.

If instead $n = p^a$ with $a \ge 1$, then $(1+X)^{p^a} = 1 + X^{p^a}$ in $\F_p[X]$, so $p \mid \binom{p^a}{j}$ for all $1 \le j \le p^a - 1$. On the other hand,
\[
\binom{p^a}{p^{a-1}} = p \binom{p^a - 1}{p^{a-1} - 1},
\]
and in $\F_p[X]$ we have $(1+X)^{p^a - 1} = (1+X^{p^a})/(1+X) = \sum_{k=0}^{p^a-1} (-1)^k X^k$, so $\binom{p^a-1}{p^{a-1}-1} \equiv (-1)^{p^{a-1}-1} \not\equiv 0 \pmod{p}$. Hence $p^2 \nmid \binom{p^a}{p^{a-1}}$, and the gcd equals $p$.
\end{proof}

Now assume $n \ge 2$ is not a prime power, and set $m := \lfloor n/2 \rfloor$. By Lemma~\ref{lem:binom-gcd} and the symmetry $\binom{n}{j} = \binom{n}{n-j}$, there exist integers $A_1, \dots, A_m$ such that
\[
1 + \sum_{j=1}^m A_j \binom{n}{j} = 0.
\]
Define integers $d_0, d_1, \dots, d_n$ by setting $d_0 = d_n = 1$, $d_j = d_{n-j} = A_j$ for $1 \le j < n/2$, and additionally $d_{n/2} = 2 A_m$ when $n$ is even. Then
\begin{equation}
\label{eq:winding-sum-zero}
\frac{1}{2} \sum_{j=0}^n \binom{n}{j} d_j = 0.
\end{equation}

We now realize the numbers $d_j$ as winding numbers of branch sums of a positively homogeneous homeomorphism $h : \C \to \C$. Set
\[
c_j := 1 - \frac{2j}{n} \quad (0 < j < n), \qquad r := \frac{1}{4n}.
\]
Notice that the punctures $(c_j, 0)$ lie in $(-1, 1) \times \{0\}$, are symmetric about the origin ($c_{n-j} = -c_j$), and are separated by distance $2/n > 2r$.

In the $(s,b)$-plane, build a piecewise-linear path $\gamma_+$ starting at $(0,r)$ as follows:
\begin{itemize}
\item For each index $j$ with $c_j > 0$ (i.e.\ $1 \le j < n/2$), travel horizontally at height $r$ from $(0,r)$ to $(c_j, r)$, traverse the diamond with center $(c_j, 0)$ and vertex radius $r$ for $d_j - 1$ signed full turns (counterclockwise for positive turns, clockwise for negative turns), and return horizontally along height $r$ back to $(0,r)$.
\item Next, traverse the diamond centered at $(0,0)$ of radius $r$ for $k$ signed full turns, where $k := A_m$ if $n$ is even and $k := 0$ if $n$ is odd.
\item Finally, move clockwise from $(0,r)$ along the right half of the central diamond via $(r,0)$ to $(0,-r)$.
\end{itemize}
Form a closed curve $\gamma$ by concatenating $\gamma_+$ with its antipodal image $-\gamma_+$ (traversed in the same parameter direction, from $(0,-r)$ back to $(0,r)$).

\begin{lemma}
\label{lem:homeo-h}
There exists a homeomorphism $h : \C \to \C$ together with a continuous isotopy $H : [0,1] \times \C \to \C$ (such that each $h_t := H(t,\cdot)$ is a homeomorphism, $h_0 = \mathrm{id}_\C$, and $h_1 = h$) satisfying:
\begin{enumerate}[label=\textup{(\alph*)}]
\item $h_t(0) = 0$ and $h_t(\rho z) = \rho h_t(z)$ for all $t \in [0,1]$, $\rho > 0$, $z \in \C$;
\item $h_t(z) - h_t(-z) = 2z$ for all $t \in [0,1]$ and $z \in \C$;
\item for every $0 \le j \le n$, the branch sum
\[
g_j(u) := j h(u) + (n-j) h(-u) \qquad (u \in S^1)
\]
is nowhere zero on $S^1$ and has winding number $\wind(g_j) = d_j$.
\end{enumerate}
\end{lemma}

\begin{proof}
The closed piecewise-linear curve $\gamma$ avoids every puncture $(c_j, 0)$ ($0 < j < n$) and satisfies $|s| < 1$ everywhere because $c_1 + r = 1 - \frac{2}{n} + \frac{1}{4n} < 1$. The horizontal outgoing and returning segments cancel in signed winding contribution, and the diamonds around distinct punctures are disjoint. Since the antipodal map $(s,b) \mapsto (-s,-b)$ is orientation-preserving in $\R^2$, the second half $-\gamma_+$ winds around $(-c_j, 0) = (c_{n-j}, 0)$ exactly $d_j - 1$ times without linking $(c_j, 0)$. For $c_j > 0$ we therefore have $\ind_{(c_j,0)}(\gamma) = d_j - 1$, and by antipodal symmetry $\ind_{(c_{n-j},0)}(\gamma) = d_{n-j} - 1$. At the origin $(0,0)$, the two halves together contribute $2k$ full turns plus one clockwise full turn from the two half-diamonds, giving $\ind_{(0,0)}(\gamma) = 2k - 1$; when $n$ is even, $c_{n/2} = 0$ and $2k - 1 = 2A_m - 1 = d_{n/2} - 1$.

Parametrizing $\gamma_+$ continuously on $[0,\pi]$ and extending antisymmetrically to $[0,2\pi]$ yields continuous odd functions $s_0, b_0 : S^1 \to \R$ (meaning $s_0(-u) = -s_0(u)$ and $b_0(-u) = -b_0(u)$) with $\max |s_0| < 1$ and $\ind_{(c_j,0)}(s_0, b_0) = d_j - 1$ for all $0 < j < n$. By standard uniform approximation on $S^1$ (for instance, Fej\'er means or polygonal smoothing, averaged with odd symmetrization $u \mapsto \frac{1}{2}(q(u)-q(-u))$), we obtain odd Lipschitz functions $s, b : S^1 \to \R$ arbitrarily close to $(s_0, b_0)$ in $C^0(S^1)$. Choosing the approximation sufficiently close keeps $\max |s| < 1$ and preserves all $n-1$ puncture avoidances and winding numbers $\ind_{(c_j,0)}(s,b) = d_j - 1$.

For a small real parameter $\varepsilon > 0$ and $t \in [0,1]$, define $h_t : \C \to \C$ by $h_t(0) = 0$ and
\begin{equation}
\label{eq:def-ht}
h_t(\rho u) := \rho u \bigl(1 + t s(u) + i t \varepsilon b(u)\bigr) \qquad (\rho > 0,\; u \in S^1).
\end{equation}
Property~(a) holds by construction. Because $s$ and $b$ are odd on $S^1$, for $z = \rho u$ we have $-z = \rho(-u)$, so
\[
h_t(-z) = -\rho u \bigl(1 - t s(u) - i t \varepsilon b(u)\bigr) = -2z + h_t(z),
\]
which proves the exact antipodal identity~(b).

To see that $h_t$ is a homeomorphism for all $t \in [0,1]$ once $\varepsilon > 0$ is sufficiently small, write $R_t(z) := (1 + t s(z/|z|)) z$ and $B(z) := b(z/|z|) z$ (with $R_t(0) = B(0) = 0$). Since $\max |s| < 1$, there exists $\alpha \in (0,1)$ such that $1 + t s(u) \ge \alpha$ for all $t \in [0,1]$ and $u \in S^1$. Each $R_t$ is a radial homeomorphism with explicit inverse $R_t^{-1}(w) = (1 + t s(w/|w|))^{-1} w$. Because $s$ and $b$ are Lipschitz on $S^1$, the positively homogeneous extensions $B$ and $R_t^{-1}$ are globally Lipschitz on $\C$, with Lipschitz constants bounded uniformly in $t \in [0,1]$ by constants $L_B$ and $L_R$ respectively. Choose $\varepsilon > 0$ so that $\varepsilon L_B L_R < 1$. Then for every $t \in [0,1]$, the map
\[
h_t \circ R_t^{-1}(w) = w + i t \varepsilon B(R_t^{-1}(w))
\]
is of the form $\mathrm{id}_\C + \psi_t$ with $\Lip(\psi_t) \le \varepsilon L_B L_R < 1$. By the Banach fixed-point theorem, $\mathrm{id}_\C + \psi_t$ is a bi-Lipschitz homeomorphism of $\C$, and its inverse varies continuously in $(t,w)$ since the contraction constant is uniform in $t$. Composing with $R_t$ shows that $h_t$ is a homeomorphism and $H(t,z) = h_t(z)$ is a continuous isotopy from $\mathrm{id}_\C$ to $h := h_1$.

Finally, let $0 < j < n$ and $u \in S^1$. Substituting~\eqref{eq:def-ht} at $t=1$ and using the oddness of $s$ and $b$ gives
\[
g_j(u) = j u(1 + s(u) + i\varepsilon b(u)) - (n-j) u(1 - s(u) - i\varepsilon b(u))
= u \Bigl( n\bigl(s(u) - c_j\bigr) + i n\varepsilon b(u) \Bigr),
\]
since $j - (n-j) = -n c_j$. Hence
\[
\frac{g_j(u)}{u} = n\bigl(s(u) - c_j\bigr) + i n\varepsilon b(u).
\]
As $(s(u), b(u))$ never hits $(c_j, 0)$, the right-hand side is nowhere zero on $S^1$. Continuously deforming the positive scaling factors $(n, n\varepsilon)$ to $(1,1)$ keeps the curve disjoint from $0$, so
\[
\wind(g_j) = \wind(u \mapsto u) + \wind\bigl(u \mapsto (s(u)-c_j) + i b(u)\bigr) = 1 + \ind_{(c_j,0)}(s,b) = d_j.
\]
For the endpoints $j=0$ and $j=n$, we have $g_n(u) = n h(u)$ and $g_0(u) = n h(-u)$. Since $h_t(S^1) \subset \C \setminus \{0\}$ gives a homotopy from $u \mapsto n u$ to $g_n$ (and from $u \mapsto -n u$ to $g_0$) in $\C \setminus \{0\}$, both $g_0$ and $g_n$ are nowhere zero and have winding number $1 = d_0 = d_n$.
\end{proof}

Set
\begin{equation}
\label{eq:def-f}
f(w) := \bigl(h^{-1}(w)\bigr)^2 \qquad (w \in \C).
\end{equation}
Since $h$ is a positively homogeneous homeomorphism of degree $1$, $f$ is continuous and positively homogeneous of degree $2$.

\begin{proof}[Proof of Proposition~\ref{prop:local-lemma}(i)]
Suppose $\sum_{i=1}^n w_i = 0$ and $f(w_1) = \dots = f(w_n) = q$. If $q = 0$, then $h^{-1}(w_i) = 0$ and hence $w_i = h(0) = 0$ for all $i$. Suppose for contradiction that $q = \zeta^2 \ne 0$ for some $\zeta \in \C \setminus \{0\}$. Then for each $i$, $h^{-1}(w_i) \in \{\zeta, -\zeta\}$, so $w_i \in \{h(\zeta), h(-\zeta)\}$. Letting $j \in \{0,\dots,n\}$ be the number of indices $i$ with $w_i = h(\zeta)$, positive homogeneity yields
\[
0 = \sum_{i=1}^n w_i = j h(\zeta) + (n-j) h(-\zeta) = |\zeta|\, g_j(\zeta/|\zeta|),
\]
contradicting Lemma~\ref{lem:homeo-h}(c).
\end{proof}

\subsection{The solution surface and finiteness of multiple readings}

Let $I := \{1, 3, 4, \dots, n\}$, and enumerate $I$ in increasing order by $j : I \to \{0, 1, \dots, n-2\}$ (so $j(1) = 0$ and $j(i) = i-2$ for $3 \le i \le n$). Define explicit positive shifts
\[
b_i := 4^{j(i)} \qquad (i \in I),
\]
and set $\Delta_i(w) := -b_i$ (constant bounded continuous functions) and $F_i(w) := f(w) - b_i$ for $i \in I$.

Consider the algebraic curve (realized as a closed subset of $\C \times \C^I$) consisting of all configurations where the $n-1$ values $F_i(w_i)$ ($i \in I$) share a common value $q \in \C$:
\[
S := \bigl\{(q, (z_i)_{i \in I}) \in \C \times \C^I : z_i^2 = q + b_i \text{ for all } i \in I\bigr\},
\]
with $w_i = h(z_i)$. Since the branch values $-b_i = -4^{j(i)}$ are pairwise distinct, at each point $(q, z) \in S$ at most one coordinate $z_i$ vanishes. Near a point with all $z_i \ne 0$, the projection $q$ provides a local holomorphic coordinate; near a branch point where $z_{i_0} = 0$, the coordinate $z_{i_0}$ provides a local holomorphic coordinate (with $q = z_{i_0}^2 - b_{i_0}$ and $z_i = \pm\sqrt{z_{i_0}^2 - b_{i_0} + b_i}$ holomorphic for $i \ne i_0$). Thus $S$ is a smooth oriented $2$-manifold (a Riemann surface). Moreover, $S$ is connected: in the unbranched covering over $\C \setminus \{-b_i : i \in I\}$, a small counterclockwise loop around $-b_i$ flips the sign of $z_i$ alone while leaving all other $z_{i'}$ unchanged, so monodromy acts transitively on the $2^{n-1}$ sheets. For any $T > 4^{n-2}$, the preimage $S_T := \{(q, z) \in S : |q| \le T\}$ is a connected compact oriented surface with boundary $\partial S_T = \{|q| = T\}$, and its boundary orientation is given by positively oriented parametrization of $|q| = T$.

Define the \emph{reading map} $e : S \to \C^2$ by
\[
e(q, z) := (L, R) := \Bigl(h(z_1),\; \sum_{i=3}^n h(z_i)\Bigr).
\]
In the notation of Proposition~\ref{prop:local-lemma}(ii), if $w_i = h(z_i)$ for $i \in I$, then $p_1 = w_1 = L$ and $p_2 = -\sum_{i=3}^n w_i = -R$ (because $p_n = 0$), so $e(q,z) = (p_1, -p_2)$ and $w_2 = p_2 - p_1 = -L - R$.

\begin{lemma}
\label{lem:e-proper-finite}
The map $e : S \to \C^2$ is proper, and its multiple-preimage locus
\[
\Sigma := \{x \in S : \exists y \in S,\; y \ne x \text{ and } e(y) = e(x)\}
\]
is finite.
\end{lemma}

\begin{proof}
Because $L = h(z_1)$ and $h$ is a homeomorphism, $z_1 = h^{-1}(L)$ is continuously determined by $L$, and therefore $q = z_1^2 - b_1$ is determined by $L$. If $(L, R)$ lies in a bounded subset of $\C^2$, then $z_1$, $q$, and every $z_i$ (via $|z_i|^2 = |q + b_i|$) are bounded. Since $S$ is closed in $\C \times \C^I$, $e$ is proper.

Now suppose $x = (q, z)$ and $y = (q', z')$ are distinct points of $S$ with $e(x) = e(y)$. Matching the first coordinate $h(z_1) = h(z'_1)$ gives $z_1 = z'_1$ and hence $q = z_1^2 - b_1 = (z'_1)^2 - b_1 = q'$. For every $i \in \{3, \dots, n\}$, we have $(z'_i)^2 = q + b_i = z_i^2$, so $z'_i \in \{z_i, -z_i\}$. Let $J := \{i \in \{3, \dots, n\} : z'_i \ne z_i\}$. Because $x \ne y$, $J$ is nonempty, and for every $i \in J$ we have $z'_i = -z_i \ne 0$. Matching the second coordinate $R$ and applying the exact antipodal identity $h(z) - h(-z) = 2z$ from Lemma~\ref{lem:homeo-h}(b), we obtain
\[
0 = \sum_{i=3}^n \bigl(h(z_i) - h(z'_i)\bigr) = \sum_{i \in J} \bigl(h(z_i) - h(-z_i)\bigr) = 2 \sum_{i \in J} z_i.
\]
Hence $\sum_{i \in J} z_i = 0$. Define the polynomial
\[
P_J(q) := \prod_{\varepsilon \in \{\pm 1\}^J} \sum_{i \in J} \varepsilon_i \sqrt{q + b_i}.
\]
Expanding the product, every term is invariant under flipping the sign of any single square root $\sqrt{q + b_i}$, so $P_J(q)$ is a polynomial in $q$ (with integer coefficients). Any $q$ for which $\sum_{i \in J} z_i = 0$ must satisfy $P_J(q) = 0$. At $q = 0$, the square roots are $\sqrt{b_i} = 2^{j(i)}$. Let $k := \max_{i \in J} j(i)$; then for every sign choice $\varepsilon \in \{\pm 1\}^J$,
\[
\left|\sum_{i \in J} \varepsilon_i 2^{j(i)}\right| \ge 2^k - \sum_{r=0}^{k-1} 2^r = 1 > 0.
\]
Therefore $P_J(0) \ne 0$, so $P_J$ is not the zero polynomial and has only finitely many roots in $\C$. Since there are only finitely many nonempty subsets $J \subset \{3,\dots,n\}$, and each $q \in \C$ has at most $2^{n-1}$ preimages in $S$, the set $\Sigma$ is finite.
\end{proof}

\subsection{Boundary winding numbers and single-cell extension}

Define the continuous residual function $G : S \to \C$ by
\[
G(q, z) := f(-L - R) - q = f\Bigl(-\sum_{i \in I} h(z_i)\Bigr) - q.
\]
Thus $G(q, z) = 0$ precisely when $w_2 = -L - R$ has the same $f$-value $q$ as the shifted values $F_i(w_i)$ ($i \in I$).

\begin{lemma}
\label{lem:boundary-winding}
There exists $T_0 > 4^{n-2}$ such that for all $T \ge T_0$:
\begin{enumerate}[label=\textup{(\alph*)}]
\item $G$ has no zeros on $\{|q| \ge T\}$, and all branch points of $S$ and all points of $\Sigma$ lie in the interior of $S_T$;
\item the sum of the winding numbers of $G$ over all boundary components of $\partial S_T$ (with the induced boundary orientation) equals $0$.
\end{enumerate}
\end{lemma}

\begin{proof}
Over the circle $|q| = T$ (with $T > 4^{n-2}$), the curve $S$ breaks into $2^{n-2}$ disjoint connected boundary components. Indeed, parametrize $|q| = T$ by $q = T u^2$ with $u \in S^1$. As $u$ makes one full counterclockwise turn around $S^1$, $q = T u^2$ winds twice around the origin, so every square root $\sqrt{T u^2 + b_i}$ returns to itself, while shifting $u \mapsto -u$ reverses all $n-1$ signs simultaneously. Consequently, each boundary component corresponds to an antipodal pair $\{\sigma, -\sigma\}$ of sign vectors $\sigma \in \{\pm 1\}^I$, and is positively parametrized by $u \in S^1$ via
\[
q(u) = T u^2, \qquad z_i(u) = \sigma_i u \sqrt{T + b_i u^{-2}} = \sqrt{T}\,\bigl(\sigma_i u + O(T^{-1})\bigr) \qquad (i \in I),
\]
where the principal square root is taken and the $O(T^{-1})$ error is uniform in $u \in S^1$.

Let $j \in \{0, \dots, n-1\}$ be the number of entries of $\sigma$ equal to $+1$. By positive homogeneity of $h$ (degree $1$) and $f$ (degree $2$), as $T \to \infty$ we have uniformly in $u \in S^1$:
\[
\frac{1}{\sqrt{T}} \Bigl(-\sum_{i \in I} h(z_i(u))\Bigr) \longrightarrow U_j(u) := -j h(u) - (n-1-j) h(-u),
\]
and hence
\[
\frac{G(q(u), z(u))}{T} \longrightarrow G^\infty_j(u) := f(U_j(u)) - u^2.
\]
We claim that $G^\infty_j(u) \ne 0$ for all $u \in S^1$. Indeed,
\[
G^\infty_j(u) = \bigl(h^{-1}(U_j(u))\bigr)^2 - u^2 = \bigl(h^{-1}(U_j(u)) - u\bigr)\bigl(h^{-1}(U_j(u)) + u\bigr).
\]
If $h^{-1}(U_j(u)) = u$, then $U_j(u) = h(u)$, so $(j+1) h(u) + (n-1-j) h(-u) = g_{j+1}(u) = 0$; if $h^{-1}(U_j(u)) = -u$, then $U_j(u) = h(-u)$, so $j h(u) + (n-j) h(-u) = g_j(u) = 0$. Both are impossible by Lemma~\ref{lem:homeo-h}(c).

By compactness of $S^1$, there exists $T_0 > 4^{n-2}$ (also strictly larger than $\max_{(q,z) \in \Sigma} |q|$ by Lemma~\ref{lem:e-proper-finite}) such that for every $T \ge T_0$, $G/T$ is uniformly close to $G^\infty_j$ on each boundary component of $\partial S_T$, and in particular $G$ is nowhere zero on $\{|q| \ge T_0\}$ and has the same winding number as $G^\infty_j$ on the component corresponding to $\{\sigma, -\sigma\}$.

To compute $\wind(G^\infty_j)$, let $h_t$ ($t \in [0,1]$) be the isotopy of Lemma~\ref{lem:homeo-h} from $\mathrm{id}_\C$ to $h$. Because $h_t$ is a homeomorphism for every $t \in [0,1]$, the homotopy
\[
t \longmapsto \bigl(h_t(h^{-1}(U_j(u))) - h_t(u)\bigr)\bigl(h_t(h^{-1}(U_j(u))) - h_t(-u)\bigr)
\]
never vanishes on $[0,1] \times S^1$. At $t = 0$ it equals $G^\infty_j(u)$; at $t = 1$ it equals
\[
\bigl(U_j(u) - h(u)\bigr)\bigl(U_j(u) - h(-u)\bigr) = \bigl(-g_{j+1}(u)\bigr)\bigl(-g_j(u)\bigr) = g_{j+1}(u)\, g_j(u).
\]
Therefore the winding number of $G$ on this boundary component is
\[
\wind(G^\infty_j) = \wind(g_{j+1}) + \wind(g_j) = d_{j+1} + d_j.
\]
Note that replacing $\sigma$ by $-\sigma$ replaces $j$ by $n-1-j$, which leaves $d_j + d_{j+1}$ unchanged because $d_{n-r} = d_r$. Summing over all $2^{n-1}$ sign vectors $\sigma \in \{\pm 1\}^I$ counts each of the $2^{n-2}$ boundary components twice, so by Pascal's identity and~\eqref{eq:winding-sum-zero}, the sum of the boundary winding numbers over $\partial S_T$ is
\[
\frac{1}{2} \sum_{j=0}^{n-1} \binom{n-1}{j} (d_j + d_{j+1}) = \frac{1}{2} \sum_{j=0}^n \binom{n}{j} d_j = 0. \qedhere
\]
\end{proof}

\begin{remark}
\label{rem:allowed-windings}
In the construction above, pairing $d_j = d_{n-j}$ ($1 \le j < n/2$) and setting $d_{n/2} = 2 A_m$ (for $n$ even) lets the total boundary winding number $\frac{1}{2}\sum_{j=0}^n \binom{n}{j} d_j$ sweep out precisely the coset $1 + \gcd_{0<j<n}\binom{n}{j}\,\Z$. Thus the topological obstruction vanishes if and only if $n$ is not a prime power. For $n = 6$, taking $(A_1, A_2, A_3) = (-1, -1, 1)$ gives $1 - 6 - 15 + 20 = 0$ and $(d_0,\dots,d_6) = (1,-1,-1,2,-1,-1,1)$; the $16$ boundary components have winding numbers $-2$ (five times), $+1$ (ten times), and $0$ (once).
\end{remark}

\begin{proof}[Proof of Proposition~\ref{prop:local-lemma}(ii)]
Fix $T > T_0$ as in Lemma~\ref{lem:boundary-winding}. We first construct a continuous nowhere-vanishing function $\widetilde G : S \to \C \setminus \{0\}$ such that $\widetilde G = G$ outside $S_T$ and $\widetilde G(x) = 1$ for every $x \in \Sigma$.

On the connected compact oriented surface $S_T$ with boundary $\partial S_T$, consider the normalized boundary map $\widehat G := G/|G| : \partial S_T \to S^1$. In the long exact cohomology sequence of the pair $(S_T, \partial S_T)$ with integer coefficients,
\[
H^1(S_T; \Z) \longrightarrow H^1(\partial S_T; \Z) \xrightarrow{\;\delta\;} H^2(S_T, \partial S_T; \Z) \cong \Z,
\]
the connecting homomorphism $\delta$ evaluates the sum of the winding numbers over the oriented components of $\partial S_T$. By Lemma~\ref{lem:boundary-winding}(b), $\delta([\widehat G]) = 0$, so $[\widehat G] \in H^1(\partial S_T; \Z)$ lies in the image of $H^1(S_T; \Z)$. Because $S^1$ is a $K(\Z, 1)$ space, this means $\widehat G$ extends to a continuous map $u_0 : S_T \to S^1$ up to homotopy on $\partial S_T$; using a collar neighborhood of $\partial S_T$ in $S_T$ to perform the homotopy, we obtain a continuous map $u : S_T \to S^1$ satisfying $u|_{\partial S_T} = \widehat G$. Extending $\log |G| : \partial S_T \to \R$ continuously to a real function $\lambda : S_T \to \R$ by Tietze's extension theorem, the function $G_1 := e^\lambda u : S_T \to \C \setminus \{0\}$ is continuous, nowhere zero, and agrees with $G$ on $\partial S_T$.

Next, because the finite set $\Sigma$ lies in $\inte S_T$, choose pairwise disjoint closed coordinate disks $B_x \subset \inte S_T$ centered at each $x \in \Sigma$. On each $B_x$, choose a complex number $\alpha_x \in \C$ with $e^{\alpha_x} = G_1(x)^{-1}$ and a continuous cutoff $\rho_x : S_T \to [0,1]$ supported in $B_x$ with $\rho_x(x) = 1$. Defining
\[
\widetilde G(y) :=
\begin{cases}
G_1(y) \exp\bigl(\sum_{x \in \Sigma} \rho_x(y) \alpha_x\bigr) & \text{if } y \in S_T, \\
G(y) & \text{if } y \in S \setminus S_T,
\end{cases}
\]
gives a continuous nowhere-vanishing function $\widetilde G : S \to \C \setminus \{0\}$ that agrees with $G$ on $S \setminus \inte S_T$ and equals $1$ at every point of $\Sigma$.

Now consider the difference $\widetilde G - G : S \to \C$. It vanishes outside the compact set $S_T$, so it is bounded and compactly supported. We check that $\widetilde G - G$ is constant on every fiber of $e : S \to \C^2$. Indeed, if $e(x) = e(y)$ with $x \ne y$, then $x, y \in \Sigma$, so $\widetilde G(x) = \widetilde G(y) = 1$; moreover, $e(x) = e(y)$ implies $L(x) = L(y)$ and $R(x) = R(y)$, which forces $q(x) = (h^{-1}(L(x)))^2 - b_1 = q(y)$ and hence $G(x) = f(-L(x)-R(x)) - q(x) = G(y)$. Therefore $\widetilde G - G$ factors through the image $e(S)$. Because $e : S \to \C^2$ is proper (Lemma~\ref{lem:e-proper-finite}) and $\C^2$ is locally compact Hausdorff, $e$ is a closed map onto the closed subset $e(S) \subset \C^2$, so the induced function $D_0 : e(S) \to \C$ satisfying $D_0 \circ e = \widetilde G - G$ is continuous, bounded, and supported in the compact set $e(S_T)$. Applying the bounded Tietze extension theorem separately to the real and imaginary parts of $D_0$ yields a bounded continuous extension $D : \C^2 \to \C$ (which may further be multiplied by a compactly supported cutoff equal to $1$ on $e(S_T)$).

Finally, let $p_0, \dots, p_n \in \C$ satisfy $p_0 = p_n = 0$ and set $w_i = p_i - p_{i-1}$. Suppose for contradiction that all $n$ values in Proposition~\ref{prop:local-lemma}(ii) equal a common number $q \in \C$. For every $i \in I$, we have $f(w_i) - b_i = q$; setting $z_i := h^{-1}(w_i)$ gives $z_i^2 = q + b_i$, so $(q, z) \in S$. By definition of $e$,
\[
e(q, z) = \Bigl(w_1,\; \sum_{i=3}^n w_i\Bigr) = (p_1, -p_2),
\]
and $w_2 = p_2 - p_1 = -L - R$. Evaluating the second value gives
\[
\bigl(f(w_2) + D(p_1, -p_2)\bigr) - q = G(q, z) + D(e(q, z)) = \widetilde G(q, z) \ne 0,
\]
a contradiction.
\end{proof}

% ============================================================================
\section{Cylinder Rigidity}
\label{sec:cylinder}
% ============================================================================

Write coordinates in $\R^d$ ($d \ge 3$) as $(u, z) \in \R^{d-1} \times \R$. We begin by proving that inside a centered cylinder, vanishing horizontal first moments force every convex partition to consist purely of horizontal slabs.

\begin{lemma}[Cylinder rigidity]
\label{lem:cylinder-rigidity}
Let $D \subset \R^{d-1}$ be a convex body whose centroid (with respect to $(d-1)$-dimensional Lebesgue measure) is $0$, and let $A := D \times [-h, h]$ for some $h > 0$. Suppose $C_1, \dots, C_n \in \Kcal(\R^d)$ are convex bodies with mutually disjoint interiors covering $A$ such that
\[
\int_{C_i} u \, dx = 0 \qquad (1 \le i \le n).
\]
Then there exist $-h = t_0 < t_1 < \dots < t_n = h$ and a permutation $\sigma \in S_n$ such that $C_{\sigma(i)} = D \times [t_{i-1}, t_i]$ for every $1 \le i \le n$.
\end{lemma}

\begin{proof}
Let $\nu$ denote the normalized $(d-1)$-dimensional Lebesgue probability measure on $D$. For each $i \in \{1,\dots,n\}$, let $\pi_i \subset D$ be the compact convex horizontal projection of $C_i$, and for $u \in D$ define the vertical fiber length
\[
\ell_i(u) := \mathcal{H}^1\bigl(\{z \in [-h,h] : (u,z) \in C_i\}\bigr),
\]
so that $\ell_i(u) = 0$ for $u \notin \pi_i$. By Fubini's theorem, $m_i := \int_D \ell_i \, d\nu = \vol_d(C_i)/\vol_{d-1}(D) > 0$, and $\sum_{i=1}^n m_i = 2h$.

Because $C_i$ is convex, the upper and lower boundary functions $z_i^+(u) := \max\{z : (u,z) \in C_i\}$ and $z_i^-(u) := \min\{z : (u,z) \in C_i\}$ are concave and convex on $\pi_i$, respectively, so $\ell_i = z_i^+ - z_i^-$ is concave on $\pi_i$. In particular, $z_i^+$ and $z_i^-$ are continuous on $\inte \pi_i$. Moreover, $\ell_i(u) > 0$ for all $u \in \inte \pi_i$ (otherwise $C_i$ would lie in a hyperplane over $\pi_i$, contradicting $\vol_d(C_i) > 0$). It follows that for every $u \in \inte \pi_i$, the open vertical segment $\{u\} \times (z_i^-(u), z_i^+(u))$ is contained in $\inte C_i$.

Consider the probability measure $\lambda_i := (\ell_i / m_i)\, d\nu$ supported on $\pi_i$. Its centroid is
\[
\int_{\pi_i} u \, d\lambda_i(u) = \frac{1}{m_i \vol_{d-1}(D)} \int_{C_i} u \, dx = 0.
\]
Because $\lambda_i$ is absolutely continuous with respect to Lebesgue measure on $\pi_i$ and $\vol_{d-1}(\pi_i) > 0$, the centroid $0$ lies in $\inte \pi_i$. Since $\ell_i$ is finite and concave on $\pi_i$, it admits a supporting affine functional $u \mapsto \ell_i(0) + L_i(u)$ on $\pi_i$ touching $\ell_i$ at $0$. Integrating $\ell_i(u) \le \ell_i(0) + L_i(u)$ against $\lambda_i$ (whose mean is $0$) and applying Cauchy--Schwarz gives
\begin{equation}
\label{eq:fiber-jensen}
\ell_i(0) \ge \int_{\pi_i} \ell_i(u) \, d\lambda_i(u) = \frac{1}{m_i} \int_D \ell_i(u)^2 \, d\nu(u) \ge \frac{1}{m_i} \left(\int_D \ell_i(u)\, d\nu(u)\right)^2 = m_i.
\end{equation}
Because $0 \in \inte \pi_i$ for every $i$, the open vertical segments $\{0\} \times (z_i^-(0), z_i^+(0))$ lie in the mutually disjoint interiors $\inte C_i$ and inside $\{0\} \times [-h, h]$. Hence
\[
2h = \sum_{i=1}^n m_i \le \sum_{i=1}^n \ell_i(0) \le 2h.
\]
Equality must hold throughout~\eqref{eq:fiber-jensen} for every $i$. Equality in Cauchy--Schwarz forces $\ell_i(u) = m_i$ for $\nu$-almost every $u \in D$. Because $\pi_i$ is closed and $\ell_i = 0$ outside $\pi_i$, we have $\pi_i = D$. Continuity of $\ell_i$ on $\inte D$ then implies $\ell_i(u) = m_i > 0$ for all $u \in \inte D$.

For every $u \in \inte D$, the $n$ open intervals $(z_i^-(u), z_i^+(u))$ are pairwise disjoint inside $[-h, h]$ and have total length $\sum_{i=1}^n m_i = 2h$, so they partition $[-h, h]$ up to endpoints. Because $z_i^-, z_i^+$ depend continuously on $u \in \inte D$ and $\inte D$ is connected, the vertical ordering of the $n$ intervals is independent of $u \in \inte D$. Relabeling by $\sigma \in S_n$ from bottom to top, we obtain $z_{\sigma(i)}^-(u) = t_{i-1}$ and $z_{\sigma(i)}^+(u) = t_i$ for all $u \in \inte D$, where $t_k := -h + \sum_{r=1}^k m_{\sigma(r)}$. Since $C_{\sigma(i)}$ is compact and full-dimensional, it is the closure of its interior, which is $\overline{\inte D \times (t_{i-1}, t_i)} = D \times [t_{i-1}, t_i]$.
\end{proof}

\begin{lemma}[Horizontal extension from an interior cylinder]
\label{lem:horizontal-extension}
Let $K \subset \R^d$ be a convex body containing $A = D \times [-h, h]$ as in Lemma~\ref{lem:cylinder-rigidity}, and let $C_1, \dots, C_n$ be a convex $n$-partition of $K$. Suppose every piece $C_i$ meets $A$ in a set of positive $d$-dimensional volume, and that the restricted pieces $C_i \cap A$ satisfy $C_i \cap A = D \times [t_{i-1}, t_i]$ for some $-h = t_0 < t_1 < \dots < t_n = h$. Then
\[
C_1 = K \cap \{z \le t_1\}, \quad C_i = K \cap \{t_{i-1} \le z \le t_i\} \;\; (2 \le i \le n-1), \quad C_n = K \cap \{z \ge t_{n-1}\}.
\]
\end{lemma}

\begin{proof}
For each $1 \le i \le n-1$, the convex bodies $C_i$ and $C_{i+1}$ have disjoint interiors, so they are separated by a hyperplane $H_i \subset \R^d$. Because $D \times [t_{i-1}, t_i] \subset C_i$ and $D \times [t_i, t_{i+1}] \subset C_{i+1}$, the hyperplane $H_i$ separates these two slabs and therefore must contain the relatively open horizontal disk $\inte D \times \{t_i\}$. Hence $H_i = \{z = t_i\}$, with $C_i \subset \{z \le t_i\}$ and $C_{i+1} \subset \{z \ge t_i\}$.

Thus $C_1 \subset K \cap \{z \le t_1\}$, $C_i \subset K \cap \{t_{i-1} \le z \le t_i\}$ for $2 \le i \le n-1$, and $C_n \subset K \cap \{z \ge t_{n-1}\}$. Because $\bigcup_{i=1}^n C_i = K$, every open slab $K \cap \{t_{i-1} < z < t_i\}$ (and the two open end slabs $\{z < t_1\}$ and $\{z > t_{n-1}\}$) must be contained in $C_i$. Taking closures and using convexity of $K$ (which implies that the closure of the intersection of $K$ with any open slab that meets $\inte K$ equals the intersection of $K$ with the corresponding closed slab) yields the stated equalities.
\end{proof}

% ============================================================================
\section{Compactness and Interior-Cylinder Localization}
\label{sec:ball-localization}
% ============================================================================

Let $K \subset \R^d$ be a convex body and $\mu$ an absolutely continuous probability measure on $\R^d$ with $\mu(K) = 1$. Let $\Pcal_\mu = \Pcal_\mu(K, n)$ denote the space of ordered equal-$\mu$-mass convex $n$-partitions $C = (C_1, \dots, C_n)$ of $K$, endowed with the product Hausdorff topology on $\Kcal(\R^d)^n$.

\begin{lemma}[Compactness of the partition space]
\label{lem:compactness}
The space $\Pcal_\mu$ is compact.
\end{lemma}

\begin{proof}
By the Blaschke selection theorem, $\{C \in \Kcal(\R^d) : C \subset K\}$ is compact, hence so is its $n$-fold product. Let $C^{(k)} = (C_1^{(k)}, \dots, C_n^{(k)}) \in \Pcal_\mu$ converge to $C = (C_1, \dots, C_n)$ in the product Hausdorff metric. For any compact convex sets $E_k \to E$ in Hausdorff distance, the indicator functions $\one_{E_k}$ converge pointwise to $\one_E$ outside $\partial E$. Because every convex set has a boundary of Lebesgue measure zero and $\mu \ll \vol_d$, dominated convergence yields $\vol_d(C_i^{(k)}) \to \vol_d(C_i)$ and $\mu(C_i^{(k)}) \to \mu(C_i) = 1/n$.

Since $\mu(C_i) = 1/n > 0$ and $\mu \ll \vol_d$, every $C_i$ has positive Lebesgue measure and is therefore full-dimensional. Moreover, the Hausdorff limit of a finite union is the union of the limits, so $\bigcup_{i=1}^n C_i = K$. Finally, $\sum_{i=1}^n \vol_d(C_i) = \lim_{k\to\infty} \sum_{i=1}^n \vol_d(C_i^{(k)}) = \vol_d(K)$. If $\inte C_i \cap \inte C_j \ne \emptyset$ for some $i \ne j$, then $\vol_d(C_i \cap C_j) > 0$, which would force $\vol_d(\bigcup_{r=1}^n C_r) < \sum_{r=1}^n \vol_d(C_r) = \vol_d(K)$, a contradiction. Hence $C \in \Pcal_\mu$.
\end{proof}

We now present the localization mechanism first on the unit Euclidean ball $K = B^d := \{x \in \R^d : |x| \le 1\}$, and then extend it in Section~\ref{sec:general-body} to arbitrary prescribed convex bodies $K$.

For $K = B^d$ and an absolutely continuous probability measure $\mu$ on $B^d$, define the cumulative height distribution
\[
F(t) := \mu\bigl(B^d \cap \{z \le t\}\bigr) \qquad (t \in [-1,1]).
\]
Because $\mu \ll \vol_d$, $F$ is continuous and nondecreasing with $F(-1) = 0$ and $F(1) = 1$. For each $1 \le i \le n-1$, the quantile level set
\[
J_i := F^{-1}(i/n) \cap [-1, 1]
\]
is a nonempty compact subinterval of $(-1, 1)$, and the intervals $J_1, \dots, J_{n-1}$ are strictly ordered from bottom to top with positive pairwise gaps (when $\mu$ is the uniform measure on $B^d$, each $J_i$ is a single point). Let $\Hcal \subset \Pcal_\mu$ be the compact set of all horizontal quantile partitions of $B^d$: partitions $(C_1, \dots, C_n)$ obtained by choosing $(b_1, \dots, b_{n-1}) \in J_1 \times \dots \times J_{n-1}$, cutting $B^d$ along the horizontal hyperplanes $z = b_i$, and permitting any permutation of the $n$ labels.

By absolute continuity of $\mu$ with respect to $\vol_d$, there exists $v_0 > 0$ such that every Borel set $E \subset B^d$ with $\mu(E) = 1/n$ satisfies $\vol_d(E) \ge v_0$. Choose $t \in (0,1)$ close enough to $1$ so that the two spherical caps $B^d \cap \{z > t\}$ and $B^d \cap \{z < -t\}$ each have volume strictly less than $v_0 / 2$. Then for every convex set $C \subset B^d$ with $\mu(C) = 1/n$, its ordinary volume centroid $(\bar u, \bar z) = \vol_d(C)^{-1} \int_C (u,z)\,dx$ satisfies the following bound. Choose $\sigma \in \{-1,1\}$ so that $\sigma \bar z = |\bar z|$ (either sign when $\bar z = 0$):
\begin{equation}
\label{eq:z-centroid-bound}
|\bar z| = \sigma \bar z \le t + (1-t)\frac{\vol_d(C \cap \{\sigma z > t\})}{\vol_d(C)} < \frac{1+t}{2} < 1,
\end{equation}
Indeed, $\sigma z \le t$ off the selected one-sided cap and $\sigma z \le 1$ on it, while $\vol_d(C \cap \{\sigma z > t\}) < \frac{1}{2}v_0 \le \frac{1}{2}\vol_d(C)$.

Fix $h \in (0,1)$ strictly larger than $(1+t)/2$ and strictly larger than $\max_{1 \le i \le n-1} \max_{b \in J_i} |b|$, and choose $a > 0$ with $a^2 + h^2 < 1$. Put
\begin{equation}
\label{eq:def-W-delta}
D := \{u \in \R^{d-1} : |u| \le a\}, \qquad A := D \times [-h, h] \subset \inte B^d, \qquad \kappa := \int_D u_1^2 \, du > 0,
\end{equation}
and for $\delta > 0$ define the continuous functional $W_\delta : \Kcal(\R^d) \to \R^{d-1}$ by
\begin{equation}
\label{eq:W-delta-formula}
W_\delta(C) := \frac{1}{\kappa} \left(\int_{C \cap A} u \, dx + \delta \int_C u \, dx\right).
\end{equation}
Because both $A$ and $B^d$ are symmetric under $(u,z) \mapsto (-u,z)$, every convex $n$-partition $\{C_i\}_{i=1}^n$ of $B^d$ satisfies $\sum_{i=1}^n W_\delta(C_i) = 0$. Continuity of $W_\delta$ on $\Kcal(\R^d)$ follows from dominated convergence, since $\partial C$ has Lebesgue measure zero.

\begin{lemma}[Localization lemma on the unit ball]
\label{lem:localization}
For every open neighborhood $U$ of $\Hcal$ in $\Pcal_\mu$, there exists $\delta_0 > 0$ such that for all $\delta \in (0, \delta_0)$, every partition $C \in \Pcal_\mu$ satisfying $W_\delta(C_i) = 0$ for all $1 \le i \le n$ belongs to $U$.
\end{lemma}

\begin{proof}
Suppose to the contrary that there exist sequences $\delta_k \downarrow 0$ and $C^{(k)} \in \Pcal_\mu \setminus U$ with $W_{\delta_k}(C_i^{(k)}) = 0$ for all $i$. By compactness of $\Pcal_\mu \setminus U$ (Lemma~\ref{lem:compactness}), pass to a convergent subsequence $C^{(k)} \to C \in \Pcal_\mu \setminus U$. Passing to the limit in~\eqref{eq:W-delta-formula} gives $\int_{C_i \cap A} u \, dx = 0$ for every $i$.

We claim that $\vol_d(C_i \cap A) > 0$ for every $i$. Suppose for contradiction that some piece $C_i$ (write $C_* := C_i$) has $\vol_d(C_* \cap A) = 0$. Then $\inte C_* \cap \inte A = \emptyset$, so there exists a nonzero vector $(p, b) \in \R^{d-1} \times \R$ weakly separating $A$ and $C_*$:
\begin{equation}
\label{eq:separation}
p \cdot u + b z \ge \max_{(u',z') \in A} (p \cdot u' + b z') = a |p| + h |b| \qquad \text{for all } (u, z) \in C_*.
\end{equation}
If $p = 0$, then $b \ne 0$ and~\eqref{eq:separation} forces $|z| \ge h$ throughout $C_*$, contradicting $|\bar z| < h$ from~\eqref{eq:z-centroid-bound}. Thus $p \ne 0$. Averaging~\eqref{eq:separation} over $C_*$ with respect to Lebesgue measure gives
\[
p \cdot \bar u \ge a |p| + h |b| - b \bar z \ge a |p| > 0,
\]
since $|b \bar z| \le |b| |\bar z| < h |b|$. Therefore $\int_{C_*} p \cdot u \, dx > 0$.

Furthermore, on $C_* \cap A$, both inequalities $p \cdot u \le a |p|$ and $b z \le h |b|$ hold, so~\eqref{eq:separation} forces $p \cdot u = a |p|$ at every point of $C_* \cap A$. By compactness and Hausdorff convergence $C_i^{(k)} \to C_*$, for all sufficiently large $k$ every point $(u, z) \in C_i^{(k)} \cap A$ satisfies $p \cdot u > a |p|/2 > 0$ (vacuously if $C_i^{(k)} \cap A = \emptyset$), and simultaneously $\int_{C_i^{(k)}} p \cdot u \, dx > 0$. Hence for large $k$,
\[
\kappa\, p \cdot W_{\delta_k}(C_i^{(k)}) = \int_{C_i^{(k)} \cap A} p \cdot u \, dx + \delta_k \int_{C_i^{(k)}} p \cdot u \, dx > 0,
\]
contradicting $W_{\delta_k}(C_i^{(k)}) = 0$.

Thus every piece $C_i$ meets $A$ in a set of positive volume and satisfies $\int_{C_i \cap A} u \, dx = 0$. By Lemma~\ref{lem:cylinder-rigidity}, the intersections $C_i \cap A$ form a horizontal slab partition of $A$, and by Lemma~\ref{lem:horizontal-extension}, $C$ is a horizontal slab partition of $B^d$. Since $\mu(C_i) = 1/n$ for all $i$, the cut heights belong to $J_1, \dots, J_{n-1}$, so $C \in \Hcal \subset U$, contradicting $C \in \Pcal_\mu \setminus U$.
\end{proof}

\begin{remark}[Direct $\delta = 0$ argument for $d = 3, n = 6$ on the uniform ball]
\label{rem:six-ball}
For the uniform probability measure on $B^3$ with $n = 6$, the global term $\delta \int_C u\,dx$ can be omitted entirely ($\delta = 0$). Taking $a = h = 1/\sqrt{2}$, the cylinder $A \subset B^3$ contains the concentric ball of radius $1/\sqrt{2}$. Any convex piece disjoint from $\inte A$ is separated from that inner ball by a hyperplane at distance at least $1/\sqrt{2}$ from the origin, and therefore lies in a spherical cap of height $1 - 1/\sqrt{2}$. The volume fraction of a cap $\{z \ge s\}$ in $B^3$ is $(2 - 3s + s^3)/4$; at $s = 1/\sqrt{2}$ this equals $\frac{1}{2} - \frac{5}{8\sqrt{2}} < \frac{243}{4000} < \frac{1}{6}$. Hence every equal-volume piece must intersect $\inte A$, and all five horizontal quantile heights lie in $(-1/\sqrt{2}, 1/\sqrt{2})$.
\end{remark}

% ============================================================================
\section{Uniform Local Charts and Global Assembly on the Unit Ball}
\label{sec:ball-assembly}
% ============================================================================

We now show that in a uniform open neighborhood $U$ of $\Hcal$ in $\Pcal_\mu$, every partition has a chain of small-slope affine interfaces, and assemble the global counterexample functionals on $B^d$.

\begin{lemma}[Uniform local interface chart]
\label{lem:local-chart}
Let $K \subset \R^d$ be a convex body containing $A = D \times [-h, h]$ in its interior as in Section~\ref{sec:ball-localization}, with all quantile intervals $J_1, \dots, J_{n-1} \subset (-h, h)$. Choose pairwise disjoint compact intervals $\widetilde J_i \subset (-h, h)$ whose interiors contain $J_i$ ($1 \le i \le n-1$), ordered strictly from bottom to top, and let $\eta > 0$ be small enough that every family of affine graphs $z = b_i + a_i \cdot u$ with $b_i \in \widetilde J_i$ and $|a_i| \le \eta$ is strictly ordered over the horizontal projection of $K$ and stays strictly inside $(-h, h)$ over $D$.

Then there exist an open neighborhood $U$ of $\Hcal$ in $\Pcal_\mu$ and pairwise disjoint compact subsets $\Dcal_1, \dots, \Dcal_n \subset \Kcal(\R^d)$ such that:
\begin{enumerate}[label=\textup{(\alph*)}]
\item for every $C \in U$, there exist a unique permutation $\sigma \in S_n$ and unique interface parameters $(a_i, b_i) \in \R^{d-1} \times \widetilde J_i$ with $|a_i| \le \eta$ ($1 \le i \le n-1$) such that
\begin{align*}
C_{\sigma(1)} &= K \cap \{z \le b_1 + a_1 \cdot u\}, \\
C_{\sigma(i)} &= K \cap \{b_{i-1} + a_{i-1} \cdot u \le z \le b_i + a_i \cdot u\} \qquad (2 \le i \le n-1), \\
C_{\sigma(n)} &= K \cap \{z \ge b_{n-1} + a_{n-1} \cdot u\},
\end{align*}
and $C_{\sigma(i)} \in \Dcal_i$ for every $1 \le i \le n$;
\item for $i = 2$, the bounding interface parameters $(a_1, b_1)$ and $(a_2, b_2)$ depend continuously on the single cell $C_{\sigma(2)} \in \Dcal_2$ alone (independently of the rest of the partition).
\end{enumerate}
\end{lemma}

\begin{proof}
For each $1 \le i \le n$, define $\Dcal_i \subset \Kcal(\R^d)$ to be the collection of all closed slabs of the form $K \cap \{b_{i-1} + a_{i-1} \cdot u \le z \le b_i + a_i \cdot u\}$ as the active parameters range independently over $\{(a_k, b_k) : b_k \in \widetilde J_k,\; |a_k| \le \eta\}$ (omitting the lower inequality when $i = 1$ and the upper inequality when $i = n$).

The parameterization from the compact parameter domain to $\Dcal_i$ is continuous: for any fixed parameter, the strict open slab $K \cap \{b_{i-1} + a_{i-1} \cdot u < z < b_i + a_i \cdot u\}$ is nonempty (since the graphs stay strictly ordered over $D \subset \inte K$), so every point of the closed slab is approximated by interior points that remain inside nearby slabs, while the closed half-space inequalities pass to Hausdorff limits. Hence each $\Dcal_i$ is compact. The sets $\Dcal_1, \dots, \Dcal_n$ are pairwise disjoint because the vertical intersection of a cell $E \in \Dcal_i$ with the line $\{u = 0\}$ is an interval whose endpoints lie in $\widetilde J_{i-1}$ and $\widetilde J_i$ (or the endpoints of $K \cap \{u = 0\}$), and the $\widetilde J_k$ are pairwise disjoint. Moreover, for $E \in \Dcal_2$, the lower and upper bounding hyperplanes contain relatively open patches of $\{z = b_1 + a_1 \cdot u\}$ and $\{z = b_2 + a_2 \cdot u\}$ inside the cylinder $\inte A$, so $(a_1, b_1)$ and $(a_2, b_2)$ are uniquely determined by the set $E$. A continuous bijection from a compact space to a Hausdorff space is a homeomorphism; hence the map $E \mapsto ((a_1, b_1), (a_2, b_2))$ on $\Dcal_2$ is continuous, proving~(b).

To prove~(a), fix a reference horizontal partition $C^\circ \in \Hcal$, relabeled so that $C^\circ_1, \dots, C^\circ_n$ are ordered from bottom to top with cut heights $b^\circ_i \in J_i$. Non-adjacent pieces $C^\circ_r, C^\circ_s$ ($|r-s| \ge 2$) are at positive Euclidean distance, so for every $C$ in a sufficiently small open neighborhood $U_{C^\circ}$ of $C^\circ$ in $\Pcal_\mu$, non-adjacent pieces $C_r, C_s$ remain disjoint. For any such $C$ and any $1 \le j \le n-1$, set $P_j := \bigcup_{r=1}^j C_r$ and $Q_j := \bigcup_{r=j+1}^n C_r$. Then $P_j \cap Q_j \subset C_j \cap C_{j+1}$. Let $H_j$ be a hyperplane weakly separating the convex bodies $C_j$ and $C_{j+1}$. For each of the two open half-spaces $H_j^-, H_j^+$ bounded by $H_j$, the intersection $K \cap H_j^\pm$ is convex and hence connected; moreover, $P_j$ and $Q_j$ are closed, cover $K \cap H_j^\pm$, and are disjoint there because $P_j \cap Q_j \subset C_j \cap C_{j+1} \subset H_j$. Since $C_j$ has an interior point in one open half-space (say $H_j^-$) and $C_{j+1}$ has an interior point in $H_j^+$, connectedness forces $K \cap H_j^- \subset P_j$ and $K \cap H_j^+ \subset Q_j$. Taking closures yields $P_j = K \cap \overline{H_j^-}$ and $Q_j = K \cap \overline{H_j^+}$.

As $C \to C^\circ$, any subsequential limit of the oriented hyperplanes $H_j$ separates $C^\circ_j$ and $C^\circ_{j+1}$ and therefore must contain $\inte D \times \{b^\circ_j\}$, so it equals the horizontal hyperplane $\{z = b^\circ_j\}$. Consequently, shrinking $U_{C^\circ}$ ensures that every $H_j$ is an affine graph $z = b_j + a_j \cdot u$ with $b_j \in \widetilde J_j$ and $|a_j| \le \eta$, and hence $C_i = P_i \cap Q_{i-1} \in \Dcal_i$ for every $i$. Covering the compact set $\Hcal$ by finitely many such neighborhoods $U_{C^\circ}$ (and all their label-permuted copies) yields the desired open neighborhood $U$.
\end{proof}

Now fix the neighborhood $U$ and disjoint compact families $\Dcal_1, \dots, \Dcal_n$ for $K = B^d$ from Lemma~\ref{lem:local-chart}, and choose $\delta \in (0, \delta_0)$ from Lemma~\ref{lem:localization}. For any interface $H_i = \{z = b_i + a_i \cdot u\}$ with $b_i \in \widetilde J_i$ and $|a_i| \le \eta$, define the interface vector $P_i \in \R^{d-1}$ by
\begin{equation}
\label{eq:interface-P}
P_i := a_i + \frac{\delta}{\kappa} \int_{B^d \cap \{z \le b_i + a_i \cdot u\}} u \, dx \qquad (1 \le i \le n-1), \qquad P_0 = P_n := 0.
\end{equation}
Notice that $P_i$ is a continuous function of $(a_i, b_i)$ alone. Because $\int_D u \, du = 0$ and $\int_D u u^\top \, du = \kappa I_{d-1}$, for any graph $z = b_i + a_i \cdot u$ staying inside $(-h, h)$ over $D$ we have
\[
\int_{A \cap \{z \le b_i + a_i \cdot u\}} u \, dx = \int_D u (b_i + a_i \cdot u + h)\, du = \kappa a_i.
\]
Since $\int_A u \, dx = 0$ and $\int_{B^d} u \, dx = 0$, subtracting consecutive prefix integrals for any partition $C \in U$ (ordered via $\sigma$ from Lemma~\ref{lem:local-chart}) gives the exact telescoping identity
\begin{equation}
\label{eq:telescoping}
W_\delta(C_{\sigma(i)}) = P_i - P_{i-1} \qquad (1 \le i \le n).
\end{equation}

\begin{remark}
\label{rem:ball-cap-direction}
On $B^d$, one can also inspect the direction of $P_i$ explicitly: writing $s = b_i / \sqrt{1 + |a_i|^2}$ and integrating along the unit normal $(-a_i, 1)/\sqrt{1 + |a_i|^2}$, the axial first moment of the spherical cap is $-\omega_{d-1} (1 - s^2)^{(d+1)/2} / (d+1) < 0$, so the integral term in~\eqref{eq:interface-P} is a positive multiple of $a_i$, making $P_i$ a strictly positive multiple of $a_i$. In particular, on $B^d$ the exact zero set of $W_\delta$ is $\Hcal$. We will not need this explicit formula in Section~\ref{sec:general-body}, because Lemma~\ref{lem:localization} already guarantees that all zero-moment partitions lie in $U$, and~\eqref{eq:telescoping} suffices inside $U$.
\end{remark}

\begin{proof}[Proof of Theorem~\ref{thm:prescribed} for $K = B^d$]
Split $\R^{d-1} = \C \times \R^{d-3}$, writing $\pi : \R^{d-1} \to \C$ for the projection onto the first two coordinates and $\pi' : \R^{d-1} \to \R^{d-3}$ for the remaining $d-3$ coordinates (empty when $d = 3$). Let $f, \Delta_i$ ($i \ne 2$), and $D_* := D$ be given by Proposition~\ref{prop:local-lemma}. Define the unperturbed continuous base test $\Phi_0 : \Kcal(\R^d) \to \C \times \R^{d-3} \cong \R^{d-1}$ by
\[
\Phi_0(C) := \bigl(f(\pi W_\delta(C)),\; \pi' W_\delta(C)\bigr).
\]
Suppose $C = (C_1, \dots, C_n) \in \Pcal_\mu$ satisfies $\Phi_0(C_1) = \dots = \Phi_0(C_n)$. Because $\sum_{i=1}^n W_\delta(C_i) = 0$, the equal values of $\pi' W_\delta(C_i)$ must all vanish, and the vectors $w_i := \pi W_\delta(C_i) \in \C$ satisfy $\sum_{i=1}^n w_i = 0$ and $f(w_1) = \dots = f(w_n)$. By Proposition~\ref{prop:local-lemma}(i), $w_i = 0$ for all $i$. Hence $W_\delta(C_i) = 0$ for all $i$, so Lemma~\ref{lem:localization} implies $C \in U$.

Therefore, on the compact complement $\Pcal_\mu \setminus U$, the continuous diameter function
\[
C \longmapsto \max_{1 \le i, j \le n} \|\Phi_0(C_i) - \Phi_0(C_j)\|
\]
is strictly positive (if $\Pcal_\mu \setminus U$ is empty, any $\eta_0 > 0$ works). Let $\eta_0 > 0$ be its minimum on $\Pcal_\mu \setminus U$.

Let $M := \max\bigl(\{\|\Delta_i\|_\infty : i \ne 2\} \cup \{\|D_*\|_\infty\}\bigr) < \infty$. For any scaling parameter $\sigma > 0$, define a continuous function $E_\sigma : \bigcup_{i=1}^n \Dcal_i \to \C$ on the pairwise disjoint compact sets $\Dcal_1, \dots, \Dcal_n$ by
\[
E_\sigma(C) :=
\begin{cases}
\sigma^2 \Delta_i\bigl(\pi W_\delta(C)/\sigma\bigr) & \text{if } C \in \Dcal_i,\; i \ne 2, \\
\sigma^2 D_*\bigl(\pi P_1(C)/\sigma,\; -\pi P_2(C)/\sigma\bigr) & \text{if } C \in \Dcal_2,
\end{cases}
\]
where $P_1(C)$ and $P_2(C)$ are evaluated via~\eqref{eq:interface-P} from the continuous single-cell interface readout on $\Dcal_2$ given by Lemma~\ref{lem:local-chart}(b). Since $\bigcup_{i=1}^n \Dcal_i$ is compact in the metric space $\Kcal(\R^d)$ and $|E_\sigma| \le M \sigma^2$ there, applying the bounded Tietze extension theorem separately to the real and imaginary parts yields a continuous extension $E_\sigma : \Kcal(\R^d) \to \C$ satisfying $\sup_{C \in \Kcal(\R^d)} |E_\sigma(C)| \le \sqrt{2} M \sigma^2$.

Choose $\sigma > 0$ small enough that $2 \sqrt{2} M \sigma^2 < \eta_0$, and define $\Phi : \Kcal(\R^d) \to \C \times \R^{d-3} \cong \R^{d-1}$ by
\[
\Phi(C) := \bigl(f(\pi W_\delta(C)) + E_\sigma(C),\; \pi' W_\delta(C)\bigr).
\]
Let $C = (C_1, \dots, C_n) \in \Pcal_\mu$ be any equal-$\mu$-mass convex $n$-partition of $B^d$.
\begin{itemize}
\item If $C \notin U$, then $\max_{i,j} \|\Phi_0(C_i) - \Phi_0(C_j)\| \ge \eta_0$, so by the triangle inequality
\[
\max_{i,j} \|\Phi(C_i) - \Phi(C_j)\| \ge \eta_0 - 2\sqrt{2} M \sigma^2 > 0.
\]
\item If $C \in U$, let $\sigma_0 \in S_n$ be the bottom-to-top ordering of Lemma~\ref{lem:local-chart}(a), so $C_{\sigma_0(i)} \in \Dcal_i$ and~\eqref{eq:telescoping} gives $\pi W_\delta(C_{\sigma_0(i)}) = \pi P_i - \pi P_{i-1}$. Dividing the first ($\C$-valued) coordinate of $\Phi(C_{\sigma_0(i)})$ by $\sigma^2$ and using the degree-$2$ positive homogeneity of $f$, we obtain the values
\[
f\Bigl(\frac{\pi W_\delta(C_{\sigma_0(2)})}{\sigma}\Bigr) + D_*\Bigl(\frac{\pi P_1}{\sigma}, -\frac{\pi P_2}{\sigma}\Bigr),
\qquad
f\Bigl(\frac{\pi W_\delta(C_{\sigma_0(i)})}{\sigma}\Bigr) + \Delta_i\Bigl(\frac{\pi W_\delta(C_{\sigma_0(i)})}{\sigma}\Bigr) \quad (i \ne 2).
\]
By Proposition~\ref{prop:local-lemma}(ii) applied to the scaled parameters $\tilde p_i := \pi P_i / \sigma$, these $n$ complex numbers cannot all be equal.
\end{itemize}
In both cases, $\Phi(C_1), \dots, \Phi(C_n)$ cannot all be equal. Taking the $d-1$ real coordinates of $\Phi$ completes the proof for $K = B^d$.
\end{proof}

\begin{remark}[Polynomial volume-moment approximation on the unit ball]
\label{rem:ball-polynomial}
Because $\Pcal_\mu$ is compact, the functionals constructed above have a strictly positive minimum diameter across all $C \in \Pcal_\mu$. On the compact set $\{C \in \Kcal(\R^d) : C \subset B^d,\; \mu(C) = 1/n\}$, ordinary polynomial volume moments $C \mapsto \int_C x^\alpha \, dx$ ($\alpha \in \Z_{\ge 0}^d$) separate points: polynomials are dense in $C(B^d)$, so two full-dimensional compact convex subsets of $B^d$ with identical volume moments have identical restricted Lebesgue measures and therefore coincide. By the real Stone--Weierstrass theorem, the first two functionals can be uniformly approximated on $\{C \subset B^d : \mu(C) = 1/n\}$ by rational-coefficient polynomials in finitely many volume moments, while the remaining $d-3$ additive coordinates can be approximated in $L^1(B^d)$ by rational polynomial densities (uniformly controlling the integral error over all $C \subset B^d$). Thus on $B^d$ one can also realize the obstruction using polynomial moments, though without an explicit degree bound on the moments.
\end{remark}

% ============================================================================
\section{Arbitrary Prescribed Convex Bodies and Robust Obstruction}
\label{sec:general-body}
% ============================================================================

We now prove Theorem~\ref{thm:prescribed} in full generality for an arbitrary prescribed convex body $K \subset \R^d$, together with Corollaries~\ref{cor:robust} and~\ref{cor:classification}. Two geometric tools replace the symmetry of the unit ball: an interior chord connecting two parallel supporting hyperplanes, and a bounded positive density $g \ge 1$ that rebalances the horizontal first moment of $K$ to zero.

\begin{lemma}[Interior chord and affine normalization]
\label{lem:interior-chord}
Let $K \subset \R^d$ be a convex body. There exists an invertible affine transformation $T : \R^d \to \R^d$ such that $K' := T(K)$ satisfies:
\begin{enumerate}[label=\textup{(\alph*)}]
\item $K' \subset \R^{d-1} \times [-1, 1]$ and both hyperplanes $\{z = -1\}$ and $\{z = 1\}$ support $K'$;
\item the open vertical segment $\{0\} \times (-1, 1)$ is contained in $\inte K'$.
\end{enumerate}
In particular, for every $h \in (0,1)$, there exists $a > 0$ such that $\{u \in \R^{d-1} : |u| \le a\} \times [-h, h] \subset \inte K'$.
\end{lemma}

\begin{proof}
Pick any point $x_0 \in \inte K$, and let $B(x_0, \rho)$ be the largest closed Euclidean ball centered at $x_0$ contained in $K$. Because $\partial B(x_0, \rho) \cap \partial K \ne \emptyset$, pick a contact point $y_+ \in \partial B(x_0, \rho) \cap \partial K$. Every hyperplane supporting $K$ at $y_+$ also supports the inscribed ball $B(x_0, \rho)$ at $y_+$, and therefore coincides with the unique tangent hyperplane $H_+$ to $\partial B(x_0, \rho)$ at $y_+$. Let $H_-$ be the opposite supporting hyperplane to $K$ parallel to $H_+$, and choose a point $y_- \in K \cap H_-$.

We claim that the open segment $(y_-, y_+)$ lies entirely in $\inte K$. Otherwise, since $K$ is convex, some point $y_* \in (y_-, y_+)$ would lie in $\partial K$. Any hyperplane $H_*$ supporting $K$ at $y_*$ must contain both $y_-$ and $y_+$ (as $y_*$ is a strict convex combination of $y_-, y_+ \in K$). In particular, $H_*$ supports $K$ at $y_+$, forcing $H_* = H_+$; this contradicts $y_- \in H_* \cap H_-$ because $H_+$ and $H_-$ are distinct parallel hyperplanes.

Now choose an invertible affine transformation $T : \R^d \to \R^d$ mapping $H_-$ to $\{z = -1\}$, $H_+$ to $\{z = 1\}$, $y_-$ to $(0,-1)$, and $y_+$ to $(0,1)$. Then $K' = T(K)$ satisfies~(a) and~(b). The final assertion follows by compactness of $\{0\} \times [-h, h] \subset \inte K'$.
\end{proof}

\begin{lemma}[Positive-weight rebalancing]
\label{lem:rebalance}
Let $K \subset \R^{d-1} \times [-1, 1]$ be a convex body with $\{0\} \times (-1, 1) \subset \inte K$. There exists a bounded measurable function $g : \R^d \to \R$ with $g \ge 1$ everywhere such that
\[
\int_K u\, g(u,z)\, du\, dz = 0.
\]
\end{lemma}

\begin{proof}
Let $M := \int_K u\, dx \in \R^{d-1}$. By Lemma~\ref{lem:interior-chord}, choose a small closed cylinder $A_0 := \{|u| \le r_0\} \times [-l_0, l_0] \subset \inte K$ with $r_0, l_0 > 0$. By rotational symmetry of $A_0$ in $u$, $\int_{A_0} u\, dx = 0$ and $\int_{A_0} u u^\top \, dx = \kappa_0 I_{d-1}$ for a constant $\kappa_0 > 0$. Define
\[
g(u, z) := 1 + \frac{\one_{A_0}(u,z)}{\kappa_0}\bigl(r_0 |M| - M \cdot u\bigr).
\]
For every $(u,z) \in A_0$, Cauchy--Schwarz gives $M \cdot u \le |M| |u| \le r_0 |M|$, so $g \ge 1$ everywhere and $g$ is bounded. Moreover,
\[
\int_K u\, g(u,z)\, dx = M + \frac{r_0 |M|}{\kappa_0} \int_{A_0} u\, dx - \frac{1}{\kappa_0} \left(\int_{A_0} u u^\top \, dx\right) M = M + 0 - M = 0. \qedhere
\]
\end{proof}

\begin{proof}[Proof of Theorem~\ref{thm:prescribed} in full generality]
Because an invertible affine map $T : \R^d \to \R^d$ induces a homeomorphism of $\Kcal(\R^d)$ and pushes forward any absolutely continuous probability measure on $K$ to an absolutely continuous probability measure on $T(K)$ (preserving convex partitions and equal masses), we may apply Lemma~\ref{lem:interior-chord} and assume without loss of generality that $K \subset \R^{d-1} \times [-1, 1]$, $\{z = \pm 1\}$ support $K$, and $\{0\} \times (-1,1) \subset \inte K$.

Fix the bounded density $g \ge 1$ on $\R^d$ from Lemma~\ref{lem:rebalance} so that $\int_K u\, g\, dx = 0$. Define the quantile intervals $J_i := F^{-1}(i/n) \cap [-1, 1] \subset (-1, 1)$ for $F(t) := \mu(K \cap \{z \le t\})$, and let $\Hcal \subset \Pcal_\mu$ be the compact set of horizontal quantile partitions of $(K, \mu)$.

As in Section~\ref{sec:ball-localization}, absolute continuity $\mu \ll \vol_d$ yields $v_0 > 0$ such that every piece $C \subset K$ with $\mu(C) = 1/n$ satisfies $\vol_d(C) \ge v_0$, and hence $\int_C g\, dx \ge v_0$ because $g \ge 1$. Since $g$ is bounded, choose $t \in (0,1)$ so close to $1$ that both end caps $K \cap \{z > t\}$ and $K \cap \{z < -t\}$ have $g\,dx$-mass strictly less than $v_0 / 2$. Then for every $C \subset K$ with $\mu(C) = 1/n$, its $g$-weighted centroid
\[
(\bar u_g, \bar z_g) := \frac{1}{\int_C g\, dx} \int_C (u, z)\, g(u, z)\, dx
\]
satisfies $|\bar z_g| < (1+t)/2 < 1$, by the exact same estimate as~\eqref{eq:z-centroid-bound}.

Choose $h \in (0,1)$ strictly larger than $(1+t)/2$ and than $\max_{1 \le i \le n-1} \max_{b \in J_i} |b|$. By Lemma~\ref{lem:interior-chord}, choose $a > 0$ so that $A := D \times [-h, h] \subset \inte K$, where $D := \{u : |u| \le a\}$. Set $\kappa := \int_D u_1^2 \, du > 0$. Define the continuous functional $W_\delta : \Kcal(\R^d) \to \R^{d-1}$ by
\[
W_\delta(C) := \frac{1}{\kappa}\left(\int_{C \cap A} u \, dx + \delta \int_C u\, g\, dx\right).
\]
By Lemma~\ref{lem:rebalance} and symmetry of $A$, we have $\sum_{i=1}^n W_\delta(C_i) = 0$ for every convex $n$-partition $\{C_i\}_{i=1}^n$ of $K$.

The localization argument of Lemma~\ref{lem:localization} carries over verbatim: if a limit piece $C_*$ has $\vol_d(C_* \cap A) = 0$, the separation inequality~\eqref{eq:separation} depends only on $A$, and averaging it against the positive measure $g\, dx$ on $C_*$ yields
\[
p \cdot \bar u_g \ge a |p| + h |b| - b \bar z_g \ge a |p| > 0
\]
(where $p = 0$ is excluded by $|\bar z_g| < h$), hence $\int_{C_*} p \cdot u\, g\, dx > 0$. As before, for large $k$ this gives $\int_{C_i^{(k)} \cap A} p \cdot u \, dx \ge 0$ and $\int_{C_i^{(k)}} p \cdot u\, g\, dx > 0$, contradicting $W_{\delta_k}(C_i^{(k)}) = 0$. Lemmas~\ref{lem:cylinder-rigidity} and~\ref{lem:horizontal-extension} then force $C \in \Hcal$.

Lemma~\ref{lem:local-chart} already applies to this $K$, producing an open neighborhood $U$ of $\Hcal$ and disjoint compact families $\Dcal_1, \dots, \Dcal_n$. For each interface $z = b_i + a_i \cdot u$, define
\[
P_i := a_i + \frac{\delta}{\kappa} \int_{K \cap \{z \le b_i + a_i \cdot u\}} u\, g\, dx \qquad (1 \le i \le n-1), \qquad P_0 = P_n := 0.
\]
Subtracting consecutive prefix integrals again gives the exact identity $W_\delta(C_{\sigma(i)}) = P_i - P_{i-1}$ for every $C \in U$. Repeating the construction of $\Phi_0$, $E_\sigma$, and $\Phi$ from Section~\ref{sec:ball-assembly} verbatim yields the desired continuous functionals on $\Kcal(\R^d)$, completing the proof of Theorem~\ref{thm:prescribed}.
\end{proof}

\begin{proof}[Proof of Corollary~\ref{cor:robust}]
Let $\varphi = (\varphi_1, \dots, \varphi_{d-1})$ be the functionals of Theorem~\ref{thm:prescribed}. Suppose for contradiction that no such $\varepsilon > 0$ exists. Then for every $m \ge 1$ there exists a convex $n$-partition $C^{(m)} = (C_1^{(m)}, \dots, C_n^{(m)})$ of $K$ such that $|\mu(C_i^{(m)}) - 1/n| < 1/m$ and $|\varphi_k(C_i^{(m)}) - \varphi_k(C_j^{(m)})| \le 1/m$ for all $k, i, j$. By Blaschke's selection theorem, pass to a Hausdorff-convergent subsequence $C^{(m)} \to C = (C_1, \dots, C_n)$. Exactly as in the proof of Lemma~\ref{lem:compactness}, dominated convergence gives $\mu(C_i) = 1/n$ and $\vol_d(C_i) > 0$, so $C$ is an equal-$\mu$-mass convex $n$-partition of $K$. By continuity of $\varphi$, $\varphi_k(C_i) = \varphi_k(C_j)$ for all $k, i, j$, contradicting Theorem~\ref{thm:prescribed}.
\end{proof}

\begin{proof}[Proof of Corollary~\ref{cor:classification}]
If $n$ is a prime power, universal fairness holds by~\cite{KHA2014,BZ2014}. If $n \ge 2$ is not a prime power, universal fairness fails by Theorem~\ref{thm:prescribed}.
\end{proof}

% ============================================================================
\section{Four-Moment Rational Polynomial Obstruction on the 3-Cube}
\label{sec:cube-3d}
% ============================================================================

We now turn to Theorem~\ref{thm:cube}, beginning in dimension $d = 3$ on the cube $K = [-1,1]^2 \times [-1,1]$ with uniform (Lebesgue) measure. Set
\[
m := \frac{2}{n}, \qquad t_i := -1 + i m \quad (0 \le i \le n), \qquad c_i := t_i - \frac{m}{2} \quad (1 \le i \le n).
\]
Write points of $\R^3$ as $(u, z) = (x_1, x_2, z)$. For any measurable set $C \subset \R^3$, denote $\avg{g}{C} := \frac{n}{8} \int_C g\, dx$, so on every equal-volume piece of $[-1,1]^3$, $\avg{g}{C}$ is the ordinary average of $g$ over $C$. Recall the four moments from~\eqref{eq:four-moments}, written as $M(C) = (W(C), V(C)) \in \R^2 \times \R^2$:
\[
W(C) := \bigl(\avg{x_1}{C},\; \avg{x_2}{C}\bigr) = \avg{u}{C}, \qquad
V(C) := \bigl(\avg{(1+x_1)z}{C},\; \avg{x_2 z}{C}\bigr).
\]

Let $X_n := \Pcal_{\mathrm{unif}}([-1,1]^3, n)$ be the compact space of ordered equal-volume convex $n$-partitions of $[-1,1]^3$. Applying Lemma~\ref{lem:cylinder-rigidity} directly to the entire cube $A = [-1,1]^2 \times [-1,1]$ shows that the only partitions $C \in X_n$ with $W(C_i) = 0$ for all $i$ are the horizontal equal-thickness slab partitions (up to permutation of labels), with interfaces $z = t_i$ ($1 \le i \le n-1$).

By the same argument as in Lemma~\ref{lem:local-chart}, for any $\delta > 0$ there is an open neighborhood $U_\delta$ of the horizontal slab partitions in $X_n$ such that every $C \in U_\delta$ admits a permutation $\sigma \in S_n$ and affine interfaces $z = b_i + a_i \cdot u$ ($1 \le i \le n-1$, with $a_0 = a_n = 0$, $b_0 = -1$, $b_n = 1$) staying inside $[-1,1]$ over $[-1,1]^2$ with $|a_i| \le 3 m \delta$ and cutting out the cells $C_{\sigma(i)}$. Because the horizontal integral of $u$ over $[-1,1]^2$ vanishes, the volume of the $i$-th layer is
\[
\int_{[-1,1]^2} \bigl((b_i + a_i \cdot u) - (b_{i-1} + a_{i-1} \cdot u)\bigr)\, du = 4 (b_i - b_{i-1}).
\]
Equality of the $n$ cell volumes (each equal to $8/n = 4m$) therefore forces $b_i = t_i$ for every $0 \le i \le n$.

\begin{lemma}[Exact four-moment identities and algebraic slope inversion]
\label{lem:cube-moments}
Let $C_i$ be the affine layer
\[
C_i = \bigl\{(u, z) \in [-1,1]^2 \times \R : t_{i-1} + a_{i-1} \cdot u \le z \le t_i + a_i \cdot u\bigr\}
\]
with $t_{i-1} + a_{i-1} \cdot u \le t_i + a_i \cdot u$ on $[-1,1]^2$. Set $p_k := a_k / (3m)$ and $w := p_i - p_{i-1}$. Then
\begin{equation}
\label{eq:cube-moment-formulas}
W(C_i) = w, \qquad
\avg{u z}{C_i} = m p_{i-1} + t_i w, \qquad
\avg{z}{C_i} = c_i + 3m\Bigl(p_{i-1} \cdot w + \frac{1}{2} |w|^2\Bigr).
\end{equation}
In particular, writing $W = (W_1, W_2)$ and $V = (V_1, V_2)$ for the four moments of $C_i$, whenever $1 + 3 W_1 > 0$ the lower slope parameter $p_{i-1}$ is uniquely recovered from $(W, V)$ by
\begin{equation}
\label{eq:cube-inversion}
(p_{i-1})_2 = \frac{V_2 - t_i W_2}{m}, \qquad
(p_{i-1})_1 = \frac{V_1 - t_i W_1 - c_i - 3m (p_{i-1})_2 W_2 - \frac{3m}{2}|W|^2}{m(1 + 3 W_1)},
\end{equation}
and the upper slope parameter is $p_i = p_{i-1} + W$.
\end{lemma}

\begin{proof}
Write $z_-(u) := t_{i-1} + 3m p_{i-1} \cdot u$ and $z_+(u) := t_i + 3m p_i \cdot u$, so that $z_+(u) - z_-(u) = m(1 + 3 w \cdot u)$ and
\[
\frac{z_+(u)^2 - z_-(u)^2}{2} = \frac{z_+(u) + z_-(u)}{2} \bigl(z_+(u) - z_-(u)\bigr)
= m \Bigl(c_i + 3m \Bigl(p_{i-1} + \frac{w}{2}\Bigr) \cdot u\Bigr) (1 + 3 w \cdot u).
\]
Let $\mathbb{E}$ denote the uniform average over $u \in [-1,1]^2$. Since $\mathbb{E}[u] = 0$, $\mathbb{E}[u u^\top] = \frac{1}{3} I_2$, and all cubic moments of $u$ vanish on $[-1,1]^2$, Fubini's theorem yields:
\begin{align*}
W(C_i) &= \frac{1}{m}\, \mathbb{E}\bigl[u \cdot m(1 + 3 w \cdot u)\bigr] = 3\, \mathbb{E}[u u^\top] w = w, \\
\avg{z}{C_i} &= \mathbb{E}\Bigl[\Bigl(c_i + 3m\Bigl(p_{i-1} + \frac{w}{2}\Bigr)\cdot u\Bigr)(1 + 3 w \cdot u)\Bigr]
= c_i + 9m\, \Bigl(p_{i-1} + \frac{w}{2}\Bigr)^\top \Bigl(\frac{1}{3} I_2\Bigr) w \\
&= c_i + 3m\Bigl(p_{i-1} \cdot w + \frac{1}{2}|w|^2\Bigr), \\
\avg{u z}{C_i} &= \mathbb{E}\Bigl[u\Bigl(c_i + 3m\Bigl(p_{i-1} + \frac{w}{2}\Bigr)\cdot u\Bigr)(1 + 3 w \cdot u)\Bigr] \\
&= 3 c_i\, \mathbb{E}[u u^\top] w + 3m\, \mathbb{E}[u u^\top]\Bigl(p_{i-1} + \frac{w}{2}\Bigr)
= m p_{i-1} + \Bigl(c_i + \frac{m}{2}\Bigr) w = m p_{i-1} + t_i w.
\end{align*}
This establishes~\eqref{eq:cube-moment-formulas}. Since $V_2 = \avg{x_2 z}{C_i} = m(p_{i-1})_2 + t_i W_2$, the first equation of~\eqref{eq:cube-inversion} follows immediately. For the second, add $\avg{z}{C_i}$ and $\avg{x_1 z}{C_i}$:
\[
V_1 = \avg{(1+x_1)z}{C_i} = c_i + 3m(p_{i-1})_1 W_1 + 3m(p_{i-1})_2 W_2 + \frac{3m}{2}|W|^2 + m(p_{i-1})_1 + t_i W_1.
\]
Grouping the terms containing $(p_{i-1})_1$ gives $m(1 + 3W_1)(p_{i-1})_1$, yielding~\eqref{eq:cube-inversion}.
\end{proof}

\begin{proof}[Proof of Theorem~\ref{thm:cube} for $d = 3$]
Let $f, \Delta_i$ ($i \ne 2$), and $D_*$ be given by Proposition~\ref{prop:local-lemma}, identifying $\R^2$ with $\C$. For each $1 \le i \le n$, choose a continuous bump function $\chi_i : \R \to [0,1]$ such that $\chi_i(v) = 1$ for $|v - c_i| \le m/8$ and $\chi_i(v) = 0$ for $|v - c_i| \ge m/4$. Since the centers $c_i$ are spaced by $m$, the functions $\chi_1, \dots, \chi_n$ have pairwise disjoint supports.

Define a globally continuous readout $B_* : \R^2 \times \R^2 \to \R^2$ on all four-moment inputs $(W, V)$ by replacing the denominator in~\eqref{eq:cube-inversion} (at $i = 2$) with a strictly positive clamp:
\begin{align*}
(B_*(W, V))_2 &:= \frac{V_2 - t_2 W_2}{m}, \\
(B_*(W, V))_1 &:= \frac{V_1 - t_2 W_1 - c_2 - 3m (B_*(W, V))_2 W_2 - \frac{3m}{2}|W|^2}{\max\bigl\{m/2,\; m(1 + 3 W_1)\bigr\}}.
\end{align*}
By~\eqref{eq:cube-moment-formulas}, we can pick $\delta \in (0, 1/12)$ sufficiently small so that for every affine layer $C_i$ with $|p_{i-1}|, |p_i| \le \delta$, we have $1 + 3 W_1 > 1/2$ (so the clamped denominator equals $m(1+3W_1)$) and $|V_1(C_i) - c_i| \le m/8$ (so $\chi_j(V_1(C_i)) = \delta_{ij}$ for all $1 \le j \le n$). In particular, on the second layer $C_2$ we have $B_*(W(C_2), V(C_2)) = p_1$ and $B_*(W(C_2), V(C_2)) + W(C_2) = p_2$.

For $\sigma > 0$, define the globally continuous map $\Psi_\sigma : \R^2 \times \R^2 \to \C$ by
\[
\begin{aligned}
\Psi_\sigma(W, V) := f(W) + \sigma^2 \Bigg[ &\sum_{i \ne 2} \chi_i(V_1)\, \Delta_i(W/\sigma) \\
&+ \chi_2(V_1)\, D_*\Bigl(\frac{B_*(W,V)}{\sigma},\; -\frac{B_*(W,V) + W}{\sigma}\Bigr) \Bigg].
\end{aligned}
\]
Because the $\chi_i$ have disjoint supports and values in $[0,1]$, $|\Psi_\sigma(W, V) - f(W)| \le M \sigma^2$ uniformly on $\R^4$, where $M = \max(\{\|\Delta_i\|_\infty : i \ne 2\} \cup \{\|D_*\|_\infty\})$.

Let $U_\delta \subset X_n$ be the open neighborhood of horizontal slab partitions on which every partition is an equal-volume affine stack with all $|p_k| \le \delta$. For any $C \in X_n \setminus U_\delta$, we cannot have $f(W(C_1)) = \dots = f(W(C_n))$: since $\sum_{i=1}^n W(C_i) = 0$, Proposition~\ref{prop:local-lemma}(i) would force $W(C_i) = 0$ for all $i$, which by Lemma~\ref{lem:cylinder-rigidity} implies $C$ is a horizontal slab partition and hence in $U_\delta$. By compactness of $X_n \setminus U_\delta$, there exists $\eta_0 > 0$ such that $\max_{i,j} |f(W(C_i)) - f(W(C_j))| \ge \eta_0$ for all $C \in X_n \setminus U_\delta$.

Fix $\sigma > 0$ with $2 M \sigma^2 < \eta_0$. Then for every $C \in X_n$:
\begin{itemize}
\item if $C \notin U_\delta$, then $\max_{i,j} |\Psi_\sigma(M(C_i)) - \Psi_\sigma(M(C_j))| \ge \eta_0 - 2 M \sigma^2 > 0$;
\item if $C \in U_\delta$, then $\chi_j(V_1(C_{\sigma_0(i)})) = \delta_{ij}$, $W(C_{\sigma_0(i)}) = p_i - p_{i-1}$, and on the second layer $B_* = p_1$ and $B_* + W = p_2$; dividing $\Psi_\sigma(M(C_{\sigma_0(i)}))$ by $\sigma^2$ reproduces the system of Proposition~\ref{prop:local-lemma}(ii) at $\tilde p_i = p_i/\sigma$, so $\Psi_\sigma(M(C_i))$ cannot be constant in $i$.
\end{itemize}
By compactness of $X_n$, there exists $\gamma_0 > 0$ such that for every $C \in X_n$,
\[
\max_{1 \le i, j \le n} \max\Bigl\{\bigl|\operatorname{Re}\Psi_\sigma(M(C_i)) - \operatorname{Re}\Psi_\sigma(M(C_j))\bigr|,\; \bigl|\operatorname{Im}\Psi_\sigma(M(C_i)) - \operatorname{Im}\Psi_\sigma(M(C_j))\bigr|\Bigr\} \ge \gamma_0.
\]
For every equal-volume cell $C_i \subset [-1,1]^3$, the four moments $M(C_i)$ lie in the compact box $K_4 := [-1,1]^2 \times [-2,2] \times [-1,1] \subset \R^4$. By the Weierstrass approximation theorem on $K_4$ (and density of $\Q$ in $\R$), there exist rational polynomials $P_n, Q_n \in \Q[X_1,X_2,X_3,X_4]$ satisfying
\[
\sup_{x \in K_4} |P_n(x) - \operatorname{Re}\Psi_\sigma(x)| < \frac{\gamma_0}{4}, \qquad
\sup_{x \in K_4} |Q_n(x) - \operatorname{Im}\Psi_\sigma(x)| < \frac{\gamma_0}{4}.
\]
Setting $\gamma_n := \gamma_0 / 2 > 0$, the triangle inequality gives
\[
\max_{1 \le i, j \le n} \max\Bigl\{\bigl|P_n(M(C_i)) - P_n(M(C_j))\bigr|,\; \bigl|Q_n(M(C_i)) - Q_n(M(C_j))\bigr|\Bigr\} > \gamma_n
\]
for every $C \in X_n$. Since $M : \Kcal(\R^3) \to \R^4$ is continuous, $P_n \circ M$ and $Q_n \circ M$ are continuous on all of $\Kcal(\R^3)$, proving Theorem~\ref{thm:cube} in dimension $d = 3$.
\end{proof}

% ============================================================================
\section{Lifting the Same Polynomial Pair to All Dimensions $d \ge 3$}
\label{sec:cube-highd}
% ============================================================================

To finish the proof of Theorem~\ref{thm:cube}, we show that for $d > 3$, vanishing of the $d-3$ extra coordinate first moments forces every equal-volume convex $n$-partition of $[-1,1]^d$ to be a strict product of a $3$-dimensional equal-volume convex partition of $[-1,1]^3$ with $[-1,1]^{d-3}$.

Write $[-1,1]^d = D_3 \times E$, where $D_3 := [-1,1]^3$ carries the active coordinates $x = (x_1, x_2, z)$ and $E := [-1,1]^{d-3}$ carries the remaining coordinates $r = (x_3, \dots, x_{d-1})$.

\begin{proposition}[Product rigidity of equal-volume cube partitions]
\label{prop:product-rigidity}
Let $d > 3$ and $n \ge 1$, and let $C_1, \dots, C_n$ be an equal-volume convex $n$-partition of $K = D_3 \times E$. Suppose that the extra first moments $\avg{x_j}{C_i}$ are independent of $i$ for every $3 \le j \le d-1$. Then there exists an equal-volume convex $n$-partition $B_1, \dots, B_n$ of $D_3$ such that
\[
C_i = B_i \times E \qquad (1 \le i \le n).
\]
\end{proposition}

\begin{proof}
Because $\sum_{i=1}^n \int_{C_i} r\, dx\, dr = \int_{D_3 \times E} r\, dx\, dr = 0$ and $\int_{C_i} r\, dx\, dr$ is independent of $i$, we have $\int_{C_i} r\, dx\, dr = 0$ for every $i = 1, \dots, n$.

For each $i$ and each $r \in E$, define the compact convex $3$-dimensional slice and its volume:
\[
B_i(r) := \{x \in D_3 : (x, r) \in C_i\}, \qquad A_i(r) := \vol_3(B_i(r)).
\]
Let $\pi_i := \{r \in E : B_i(r) \ne \emptyset\}$ be the compact convex projection of $C_i$ onto $E$, so $A_i(r) = 0$ outside $\pi_i$. Let $\nu$ denote the normalized $(d-3)$-dimensional Lebesgue probability measure on $E$, and set $M := 8/n$. By Fubini's theorem,
\[
\int_E A_i(r)\, d\nu(r) = \frac{\vol_d(C_i)}{2^{d-3}} = \frac{2^d / n}{2^{d-3}} = \frac{8}{n} = M > 0.
\]

\emph{Step 1: Constant slice volume on $\inte E$.}
For any $r_0, r_1 \in \pi_i$ and $\lambda \in [0,1]$, convexity of $C_i$ gives
\[
(1-\lambda) B_i(r_0) + \lambda B_i(r_1) \subset B_i\bigl((1-\lambda) r_0 + \lambda r_1\bigr).
\]
By the $3$-dimensional Brunn--Minkowski inequality, the function $g_i(r) := A_i(r)^{1/3}$ is concave on $\pi_i$. Consider the probability measure $\lambda_i := (A_i / M)\, d\nu$ on $\pi_i$. Its centroid is $\frac{1}{M \cdot 2^{d-3}} \int_{C_i} r\, dx\, dr = 0$; since $\vol_{d-3}(\pi_i) > 0$ and $\lambda_i \ll \nu$, the point $0$ lies in $\inte \pi_i$.

Because $g_i$ is finite and concave on $\pi_i$, it admits a supporting affine majorant $r \mapsto g_i(0) + L_i(r)$ on $\pi_i$ touching $g_i$ at $0$. Integrating against $\lambda_i$ (which has mean $0$) and applying Jensen's inequality to the strictly convex function $t \mapsto t^{4/3}$ on $[0, \infty)$ yields
\begin{equation}
\label{eq:slice-jensen}
\begin{aligned}
A_i(0)^{1/3} = g_i(0) &\ge \int_{\pi_i} A_i(r)^{1/3}\, d\lambda_i(r) \\
&= \frac{1}{M} \int_E A_i(r)^{4/3}\, d\nu(r) \\
&\ge \frac{1}{M} \left(\int_E A_i(r)\, d\nu(r)\right)^{4/3} \\
&= M^{1/3}.
\end{aligned}
\end{equation}
We claim that for every $r \in \inte \pi_i$, the $3$-dimensional relative interior $\inte_3 B_i(r)$ satisfies $\inte_3 B_i(r) \times \{r\} \subset \inte C_i$. Indeed, if $x \in \inte_3 B_i(r)$ and $(x, r) \in \partial C_i$, let $(p, q) \in \R^3 \times \R^{d-3}$ be a nonzero outer normal of a supporting hyperplane to $C_i$ at $(x, r)$. Since $x \in \inte_3 B_i(r)$, varying $x'$ in a small $3$-ball around $x$ inside $B_i(r)$ gives $p \cdot (x' - x) \le 0$, forcing $p = 0$. Then $q \ne 0$ and $q \cdot (r' - r) \le 0$ for all $r' \in \pi_i$, contradicting $r \in \inte \pi_i$.

Because $0 \in \inte \pi_i$ for every $i$, the open sets $\inte_3 B_i(0)$ are pairwise disjoint subsets of $D_3$. Hence
\[
n M = 8 = \vol_3(D_3) \ge \sum_{i=1}^n A_i(0) \ge n M.
\]
Equality therefore holds throughout~\eqref{eq:slice-jensen} for every $i$. By strict convexity of $t \mapsto t^{4/3}$, we obtain $A_i(r) = M$ for $\nu$-almost every $r \in E$. Since $\pi_i$ is closed and $A_i = 0$ outside $\pi_i$, we must have $\pi_i = E$. Because the finite concave function $g_i = A_i^{1/3}$ is continuous on $\inte E$, it follows that $A_i(r) = M$ for \emph{every} $r \in \inte E$ and every $1 \le i \le n$.

\emph{Step 2: Identical slices via Brunn--Minkowski equality and second moments.}
Fix any $r \in \inte E$. Then $-r \in \inte E$ as well. For each $i$, convexity of $C_i$ gives
\[
\frac{1}{2}\bigl(B_i(r) + B_i(-r)\bigr) \subset B_i(0),
\]
and all three compact convex bodies $B_i(r), B_i(-r), B_i(0) \subset D_3$ have the same positive volume $M$. Hence equality holds in the Brunn--Minkowski inequality:
\[
\vol_3\Bigl(\frac{1}{2} B_i(r) + \frac{1}{2} B_i(-r)\Bigr)^{1/3} = \frac{1}{2} \vol_3(B_i(r))^{1/3} + \frac{1}{2} \vol_3(B_i(-r))^{1/3} = M^{1/3}.
\]
By the equality condition of the Brunn--Minkowski inequality for convex bodies of positive volume (see e.g.~\cite[Theorem~7.1.1]{Schneider2014}), $B_i(r)$ and $B_i(-r)$ are homothetic, and since they have the same volume $M > 0$, they are translates of each other. Thus there exists $v_i \in \R^3$ such that $B_i(r) = \frac{1}{2}(B_i(r)+B_i(-r)) + v_i$ and $B_i(-r) = \frac{1}{2}(B_i(r)+B_i(-r)) - v_i$. Moreover, $\frac{1}{2}(B_i(r) + B_i(-r))$ is a compact convex body contained in $B_i(0)$ with the exact same volume $M$; since a proper inclusion of full-dimensional compact convex bodies strictly increases volume, we have $\frac{1}{2}(B_i(r) + B_i(-r)) = B_i(0)$, and therefore
\[
B_i(r) = B_i(0) + v_i, \qquad B_i(-r) = B_i(0) - v_i \qquad (1 \le i \le n).
\]

At each of the three interior parameters $r, -r, 0 \in \inte E$, the slices $\{B_i(s)\}_{i=1}^n$ have mutually disjoint $3$-dimensional interiors (as shown in Step~1) and cover $D_3$ (since $\bigcup_i C_i = D_3 \times E$). Thus their total second moment over $D_3$ is independent of $s \in \{r, -r, 0\}$:
\[
\sum_{i=1}^n \int_{B_i(s)} |x|^2 \, dx = \int_{D_3} |x|^2 \, dx.
\]
Using the parallelogram identity $|x + v_i|^2 + |x - v_i|^2 - 2|x|^2 = 2|v_i|^2$, we compute
\begin{align*}
0 &= \sum_{i=1}^n \left(\int_{B_i(r)} |x|^2 \, dx + \int_{B_i(-r)} |x|^2 \, dx - 2 \int_{B_i(0)} |x|^2 \, dx\right) \\
  &= \sum_{i=1}^n \int_{B_i(0)} \bigl(|x + v_i|^2 + |x - v_i|^2 - 2|x|^2\bigr)\, dx
   = 2 M \sum_{i=1}^n |v_i|^2.
\end{align*}
Because $M > 0$, we conclude that $v_i = 0$ for every $i$, and hence $B_i(r) = B_i(0)$ for all $r \in \inte E$.

\emph{Step 3: Closure to $\partial E$.}
Set $B_i := B_i(0)$. Since $B_i \times \inte E \subset C_i$ and $C_i$ is closed, we have $B_i \times E \subset C_i$. Conversely, let $(x, r) \in C_i$ with $r \in \partial E$. For any $x_0 \in \inte_3 B_i$, we have $(x_0, 0) \in C_i$, so for every $\lambda \in (0, 1)$, convexity gives $((1-\lambda) x + \lambda x_0, (1-\lambda) r) \in C_i$ with $(1-\lambda) r \in \inte E$. Hence $(1-\lambda) x + \lambda x_0 \in B_i((1-\lambda)r) = B_i$. Letting $\lambda \downarrow 0$ and using closedness of $B_i$ yields $x \in B_i$. Therefore $C_i = B_i \times E$ for all $1 \le i \le n$, and $\{B_i\}_{i=1}^n$ is an equal-volume convex $n$-partition of $D_3$.
\end{proof}

\begin{proof}[Proof of Theorem~\ref{thm:cube} for all $d \ge 3$]
Let $P_n, Q_n \in \Q[X_1,X_2,X_3,X_4]$ be the rational polynomials constructed for $d = 3$ in Section~\ref{sec:cube-3d}. The case $d = 3$ is already proved. For $d > 3$, suppose for contradiction that $C_1, \dots, C_n$ is an equal-volume convex $n$-partition of $[-1,1]^d$ on which $P_n(M(C_i))$, $Q_n(M(C_i))$, and $\avg{x_j}{C_i}$ ($3 \le j \le d-1$) are all independent of $i$.

By Proposition~\ref{prop:product-rigidity}, there exists an equal-volume convex $n$-partition $B_1, \dots, B_n$ of $[-1,1]^3$ such that $C_i = B_i \times [-1,1]^{d-3}$ for every $1 \le i \le n$. For each of the four active integrands $g \in \{x_1, x_2, (1+x_1)z, x_2 z\}$, Fubini's theorem gives
\[
\frac{n}{2^d} \int_{B_i \times [-1,1]^{d-3}} g(x)\, dx\, dr = \frac{n}{2^d} \cdot 2^{d-3} \int_{B_i} g(x)\, dx = \frac{n}{8} \int_{B_i} g(x)\, dx.
\]
Thus $M(C_i) = M(B_i)$ for every $i$. Consequently, $P_n(M(B_i))$ and $Q_n(M(B_i))$ are independent of $i$, contradicting the $d = 3$ case of Theorem~\ref{thm:cube}.
\end{proof}

\begin{thebibliography}{9}

\bibitem{Akopyan2026}
A.~Akopyan,
\emph{Sixfold Tverberg counterexamples and convex equipartition obstructions},
manuscript, 2026,
\url{https://github.com/akopjan/toptverbert-nandakumar/blob/d98f6f87bb60b9e5b3a57370f393ade1397189bc/paper/manuscript.tex}.

\bibitem{AkopyanRepo}
A.~Akopyan,
\emph{toptverbert-nandakumar} (Lean formalization repository, commit \texttt{d98f6f87}),
2026,
\url{https://github.com/akopjan/toptverbert-nandakumar}.

\bibitem{AAK2018}
A.~Akopyan, S.~Avvakumov, and R.~Karasev,
\emph{Convex fair partitions into an arbitrary number of pieces},
Adv. Math. \textbf{493} (2026), 110927,
\url{https://doi.org/10.1016/j.aim.2026.110927};
arXiv:1804.03057v14.

\bibitem{BKS2023}
P.~V.~M. Blagojevi\'c, R.~Karasev, and P.~Sober\'on,
\emph{Equipartitions of measures by convex sets},
arXiv:2305.10711v2, 2023,
\url{https://arxiv.org/abs/2305.10711}.

\bibitem{BZ2014}
P.~V.~M. Blagojevi\'c and G.~M. Ziegler,
\emph{Convex equipartitions via equivariant obstruction theory},
Israel J. Math. \textbf{200} (2014), 49--77.

\bibitem{KHA2014}
R.~Karasev, A.~Hubard, and B.~Aronov,
\emph{Convex equipartitions: the spicy chicken theorem},
Geom. Dedicata \textbf{170} (2014), 263--279, arXiv:1306.2741.

\bibitem{NRR2012}
R.~Nandakumar and N.~Ramana Rao,
\emph{Fair partitions of polygons: an elementary introduction},
Proc. Indian Acad. Sci. Math. Sci. \textbf{122} (2012), 459--467.

\bibitem{Schneider2014}
R.~Schneider,
\emph{Convex Bodies: The Brunn--Minkowski Theory},
2nd ed., Encyclopedia of Mathematics and Its Applications, vol.~151,
Cambridge University Press, Cambridge, 2014.

\end{thebibliography}

\end{document}
